%% ==================================================
%% Chapter: Local quantum circuits
%% Sources: arXiv:2307.01696 (Malz, Styliaris, Wei, Cirac), main text before
%% Theorem 1, the caption of Fig. 1, the paragraphs "The sequential-RG
%% circuit", "Tree-RG circuit with measurements", "Preparations using
%% measurements" and "Long-range MPS using measurements", and the proof of
%% Theorem 1 in the Supplemental Material; arXiv:2103.13367 (Piroli,
%% Styliaris, Cirac), paragraphs "Quantum circuits and LOCC" and "State
%% transformations with QC and LOCC", Definition "Transformations under QC
%% and LOCC", and Examples 1 and 2.
%% ==================================================

\chapter{Local Quantum Circuits}
\label{ch:local_quantum_circuits}

A chain of $N$ sites of local dimension $d$ carries $(\C^d)^{\otimes N}$, with
the computational basis labelled by the configurations of the chain. A local
circuit is a product of layers of unitaries on disjoint pairs of neighbouring
sites of the chain closed into a ring \cite{Malz2024Preparation}, or more
generally on disjoint bonds of a geometry such as the edges of a lattice
\cite{Piroli2021LOCC}. A circuit of depth $T$ enlarges the set of sites on
which an operator acts by at most $T$ steps along the bonds, on the ring by at
most $T$ sites on each side, so a state that it prepares from a product state
has no connected correlations beyond distance $2T$, and every unitary on a bounded
number of sites is a product of a bounded number of gates on neighbouring
sites. Measurements in the computational basis, followed by corrections on
single sites chosen from all the outcomes \cite{Piroli2021LOCC}, prepare
GHZ-type states in constant depth, teleport sites along chains of neighbouring
pairs, and apply gates on pairs of distant sites or of distant registers in
constant depth. Circuits of channels on neighbouring pairs, with ancillas
attached and discarded site by site, obey the same light cone.
Chapter~\ref{ch:log_depth_preparation} applies these results to the
preparation of matrix product states.

\section{Operators on sets of sites and local circuits}
\label{sec:qc_local_circuits}

\begin{definition}[Local circuit of depth $T$]
    \label{def:qc_local_circuit}
    \lean{QuantumCircuit.IsLocalCircuitOfDepth, QuantumCircuit.IsPreparedInDepth,
        QuantumCircuit.Layer, QuantumCircuit.circuitOp,
        QuantumCircuit.BondLayer, QuantumCircuit.IsBondCircuitOfDepth,
        QuantumCircuit.IsBondPreparedInDepth, QuantumCircuit.ringBond,
        QuantumCircuit.supportedOperators,
        QuantumCircuit.rectKronecker_mem_supportedOperators}
    \leanok
    A \emph{local circuit of depth $T$} on a chain of $N$ sites is a product of
    $T$ layers, each a tensor product of unitaries acting on disjoint pairs of
    neighbouring sites. A state is \emph{prepared in depth $T$} if it equals a
    local circuit of depth $T$ applied to a product state.

    More generally, a \emph{bond geometry} on a finite set of sites $I$ is a
    family $(B_b)_{b\in\beta}$ of subsets of $I$, the bonds, and a layer is a
    product of unitaries acting on pairwise disjoint bonds. The model of
    \cite{Piroli2021LOCC} on a lattice takes for bonds the pairs of
    nearest-neighbour sites, that is the edges of a graph. Local circuits of
    depth $T$ and states prepared in depth $T$ are defined in the same way for
    every bond geometry.

    The chain is a ring with sites $1,\dots,N$: the neighbouring pairs are
    $\{k,k+1\}$ with indices modulo $N$; for $N=1$ the only pair is $\{1\}$,
    and for $N=2$ both pairs are $\{1,2\}$. The open chain has the pairs
    $\{k,k+1\}$ with $k<N$. An operator
    \emph{acts on} a set of sites $X$ if it
    lies in $\MN{d}^{\otimes X}\otimes\mathbb 1$, the span of the products
    $\bigotimes_i m_i$ with $m_i=\mathbb 1$ for $i\notin X$.

    Acting on a set of sites needs no geometry: for any finite set of sites
    $I$ and $X\subseteq I$, an operator on $(\C^d)^{\otimes I}$ acts on $X$ if
    it lies in the span of the products $\bigotimes_{i\in I}m_i$ with
    $m_i=\mathbb 1$ for $i\notin X$. The chain of $N$ sites is the case
    $I=\{1,\dots,N\}$.
\end{definition}

\begin{definition}[Pairing local coordinates]
    \label{def:qc_site_pair_coordinates}
    \lean{QuantumCircuit.sitePairConfigurationEquiv}
    \leanok
    Let $I$ be a finite set of sites and let $d,e\geq0$ be integers. At each
    site identify $(a,b)\in\{0,\dots,d-1\}\times\{0,\dots,e-1\}$ with
    $ea+b\in\{0,\dots,de-1\}$. Applied at every site, this gives a bijection
    $p$ from pairs of configurations to configurations of local dimension
    $de$.
\end{definition}

\begin{theorem}[Support of a tensor product on paired sites]
    \label{thm:qc_site_pair_tensor_support}
    \lean{QuantumCircuit.reindex_kronecker_mem_supportedOperators}
    \leanok
    \uses{def:qc_local_circuit,def:qc_site_pair_coordinates}
    Let $S,T\subseteq I$. If $A$ on $(\C^d)^{\otimes I}$ is supported on
    $S$ and $B$ on $(\C^e)^{\otimes I}$ is supported on $T$, then the matrix
    $A\otimes B$, expressed in the coordinates $p$, is supported on
    $S\cup T$ in $(\C^{de})^{\otimes I}$. This is the support identification
    for the doubled local interactions in \cite{OpenAI2026PolynomialPEPS},
    \texttt{02-information.tex}, lines 416--424 and 513--524.
\end{theorem}

\begin{proof}\leanok
    By bilinearity, it suffices to consider
    $A=\bigotimes_i A_i$ and $B=\bigotimes_i B_i$, with $A_i=I$ off $S$
    and $B_i=I$ off $T$. Pairing the local coordinates identifies their
    tensor product with $\bigotimes_i(A_i\otimes B_i)$. Every factor outside
    $S\cup T$ is the identity.
\end{proof}

\begin{definition}[Light cone of a bond geometry]
    \label{def:qc_light_cone_geometry}
    \uses{def:qc_local_circuit}
    \lean{QuantumCircuit.bondNeighbourhood, QuantumCircuit.lightCone,
        QuantumCircuit.neighbourhood, QuantumCircuit.IsSeparatedBy,
        QuantumCircuit.openBond, QuantumCircuit.edgeBond}
    \leanok
    For a bond geometry $(B_b)_{b\in\beta}$ on $I$ and $X\subseteq I$, the
    \emph{one-step neighbourhood} of $X$ is $X$ together with every bond that
    meets $X$, and the \emph{light cone of radius $r$} of $X$ is its $r$-fold
    one-step neighbourhood. On the ring, the sites within ring distance $r$ of
    $X$ are the sites $i+m$ with $i\in X$ and $\lvert m\rvert\le r$, indices
    modulo $N$, and two sets of sites are at distance larger than $s$ if no
    site of one is within ring distance $s$ of a site of the other. For a
    graph, the bonds are its edges, and the distance of two vertices is the
    least number of edges of a path joining them, as in
    \cite{Piroli2021LOCC}.
\end{definition}

\begin{lemma}[Light cones of the ring, the open chain and a graph]
    \label{lem:qc_light_cone_geometries}
    \uses{def:qc_light_cone_geometry}
    \lean{QuantumCircuit.lightCone_ringBond,
        QuantumCircuit.bondNeighbourhood_ringBond, QuantumCircuit.neighbourhood_succ,
        QuantumCircuit.disjoint_lightCone_ringBond_of_isSeparatedBy,
        QuantumCircuit.lightCone_openBond_subset, QuantumCircuit.disjoint_lightCone_openBond,
        QuantumCircuit.lightCone_edgeBond_subset,
        QuantumCircuit.lightCone_edgeBond_subset_dist,
        QuantumCircuit.disjoint_lightCone_edgeBond,
        QuantumCircuit.disjoint_lightCone_edgeBond_of_dist,
        QuantumCircuit.lightCone_lightCone}
    \leanok
    On the ring of $N$ sites, for every $N\ge1$, the light cone of radius $r$
    of $X$ is the set of sites within ring distance $r$ of $X$. On the open
    chain it is contained in the set of sites within distance $r$ of $X$ along
    the chain, and for the edges of a graph it is contained in the set of
    vertices at graph distance at most $r$ from $X$. In each case, sets of
    sites at distance larger than $2r$ have disjoint light cones of radius
    $r$. For every bond geometry, the light cone of radius $s$ of the light
    cone of radius $r$ of $X$ is the light cone of radius $r+s$ of $X$.
\end{lemma}
\begin{proof}
    \leanok
    On the ring, the bonds meeting $X$ are $\{i-1,i\}$ and $\{i,i+1\}$ for
    $i\in X$, so the one-step neighbourhood is the set of sites within ring
    distance $1$, and a displacement of modulus at most $r+1$ is one of
    modulus at most $r$ followed by one of modulus at most $1$. In the other
    two cases, two sites on a common bond are at distance at most $1$, and the
    inclusion follows by induction on $r$ and the triangle inequality. A site
    in both light cones would give, by the triangle inequality, two points of
    the two sets at distance at most $2r$. The composition law holds because
    the light cone of radius $r$ is the $r$-fold iterate of the one-step
    neighbourhood.
\end{proof}

\begin{lemma}[Light cone of a local circuit]
    \label{lem:qc_light_cone}
    \uses{def:qc_local_circuit, def:qc_light_cone_geometry,
        lem:qc_light_cone_geometries}
    \lean{QuantumCircuit.conj_circuitOp_mem_supportedOperators,
        QuantumCircuit.IsBondCircuitOfDepth.conj_mem_supportedOperators,
        QuantumCircuit.conj_circuitOp_mem_supportedOperators_lightCone,
        QuantumCircuit.BondLayer.conj_op_mem_supportedOperators}
    \leanok
    Let $U$ be a local circuit of depth $T$ on a bond geometry and let $Z$ act
    on a set of sites $X$. Then $U^\dagger ZU$ acts on the light cone of
    radius $T$ of $X$; on the ring, on the sites within ring distance $T$ of
    $X$. This is the finite light cone of \cite[Theorem~1]{Malz2024Preparation}.
\end{lemma}
\begin{proof}
    \leanok
    By induction on the layers it suffices to treat one layer $W$. Split its
    gates into those whose bond meets $X$, with product $G$, and the others,
    with product $H$. The gates commute, so $W=HG$. The operator $H$ is unitary
    and acts on the complement of $X$, so it commutes with $Z$ and
    $W^\dagger ZW=G^\dagger ZG$. The bonds of $G$ lie in the one-step
    neighbourhood of $X$, and products and adjoints of operators acting on a
    set act on it. On the ring, Lemma~\ref{lem:qc_light_cone_geometries}
    identifies the light cone.
\end{proof}

\begin{lemma}[No connected correlations beyond the light cone]
    \label{lem:qc_light_cone_correlations}
    \uses{lem:qc_light_cone}
    \lean{QuantumCircuit.expect_mul_eq_of_isPreparedInDepth,
        QuantumCircuit.expect_mul_mul_expect_one_of_isPreparedInDepth,
        QuantumCircuit.IsBondPreparedInDepth.expect_mul_eq,
        QuantumCircuit.IsBondPreparedInDepth.expect_mul_mul_expect_one,
        QuantumCircuit.expect_productVector_rectKronecker,
        QuantumCircuit.expect_productVector_mul}
    \leanok
    Let $\ket{\psi}$ be a unit vector prepared in depth $T$ on a bond
    geometry, and let $Z$, $Z'$ act on sets of sites whose light cones of
    radius $T$ are disjoint, for instance on the ring sets at distance larger
    than $2T$. Then
    $\bra{\psi}ZZ'\ket{\psi}=\bra{\psi}Z\ket{\psi}\bra{\psi}Z'\ket{\psi}$.
    Without the normalization, every vector $\ket{\psi}$ prepared in depth $T$
    satisfies
    $\bra{\psi}ZZ'\ket{\psi}\braket{\psi}{\psi}
        =\bra{\psi}Z\ket{\psi}\bra{\psi}Z'\ket{\psi}$.
    This is the statement in the proof of
    \cite[Theorem~1]{Malz2024Preparation} that every connected correlation of
    such operators vanishes. Its two ingredients on product vectors hold for
    any finite set of sites: the expectation of $\bigotimes_im_i$ in
    $\ket\pi=\bigotimes_i\ket{v_i}$ is $\prod_i\bra{v_i}m_i\ket{v_i}$, and
    $\bra{\pi}AB\ket{\pi}\braket{\pi}{\pi}
        =\bra{\pi}A\ket{\pi}\bra{\pi}B\ket{\pi}$ for $A$, $B$ acting on
    disjoint sets of sites.
\end{lemma}
\begin{proof}
    \leanok
    Write $\ket{\psi}=U\ket{\pi}$ with $\ket{\pi}=\bigotimes_i\ket{v_i}$. Since
    $UU^\dagger=\mathbb 1$,
    $\bra{\psi}ZZ'\ket{\psi}=\bra{\pi}(U^\dagger ZU)(U^\dagger Z'U)\ket{\pi}$,
    and by Lemma~\ref{lem:qc_light_cone} the two factors act on disjoint sets
    of sites. For products $\bigotimes_i m_i$ and $\bigotimes_i m'_i$ with
    disjoint supports the expectation in $\ket{\pi}$ is
    $\prod_i\bra{v_i}m_im'_i\ket{v_i}$, and at each site one of $m_i$, $m'_i$
    is the identity; bilinearity then gives
    $\bra{\pi}AB\ket{\pi}\braket{\pi}{\pi}
        =\bra{\pi}A\ket{\pi}\bra{\pi}B\ket{\pi}$ for $A$, $B$ acting on
    disjoint sets.
\end{proof}

\begin{lemma}[Inverses and series composition of local circuits]
    \label{lem:qc_circuit_inverse_composition}
    \uses{def:qc_local_circuit}
    \lean{QuantumCircuit.BondLayer.adjoint, QuantumCircuit.BondLayer.adjoint_partialOp,
        QuantumCircuit.BondLayer.adjoint_op, QuantumCircuit.circuitOp_append,
        QuantumCircuit.circuitOp_map_adjoint_reverse,
        QuantumCircuit.IsBondCircuitOfDepth.mem_unitary,
        QuantumCircuit.IsBondCircuitOfDepth.star,
        QuantumCircuit.IsBondCircuitOfDepth.mul}
    \leanok
    For a layer $W$ with gates $g_k$, the layer with the gates $g_k^\dagger$
    on the same bonds implements $W^\dagger$, and the same holds for the
    product of the gates of $W$ over any subset of its bonds. If a circuit
    has the layers $W_1,\dots,W_T$, applied in this order, and implements $U$,
    then the circuit with the layers $W_T^\dagger,\dots,W_1^\dagger$
    implements $U^\dagger$; the layers of a circuit implementing $U$ followed
    by those of a circuit implementing $U'$ implement $U'U$. Consequently, a
    local circuit of depth $T$ is unitary, and its adjoint is a local circuit
    of depth $T$. If $U$ and $U'$ are local circuits of depths $T$ and $T'$,
    then $U'U$ is a local circuit of depth $T+T'$.
\end{lemma}
\begin{proof}
    \leanok
    If $U=W_T\cdots W_1$ with layers $W_k$, then
    $U^\dagger=W_1^\dagger\cdots W_T^\dagger$, and $W_k^\dagger$ is the layer
    whose gates are the adjoints of the gates of $W_k$, on the same bonds. The
    layers of $U'U$ are those of $U$ followed by those of $U'$.
\end{proof}

\begin{lemma}[Preparation from the all-zero state]
    \label{lem:qc_preparation_from_zero}
    \uses{def:qc_local_circuit}
    \lean{QuantumCircuit.rectKronecker_mulVec_productVector,
        QuantumCircuit.finKronecker_mulVec_productVector,
        Matrix.exists_mem_unitary_eq_smul_mulVec_single_zero,
        QuantumCircuit.productVector_pi_smul,
        QuantumCircuit.exists_productVector_eq_smul_finKronecker_mulVec,
        QuantumCircuit.isLocalCircuitOfDepth_finKronecker,
        QuantumCircuit.IsPreparedInDepth.exists_eq_smul_mulVec_productVector_single_zero}
    \leanok
    Let $d\ge1$. For one-site operators $m_i$ and vectors $\ket{v_i}$ in
    $\mathbb{C}^d$ on any finite set of sites,
    $(\bigotimes_im_i)(\bigotimes_i\ket{v_i})=\bigotimes_im_i\ket{v_i}$, and
    $\bigotimes_ic_i\ket{v_i}=(\prod_ic_i)\bigotimes_i\ket{v_i}$ for scalars
    $c_i$. Every nonzero $x\in\mathbb{C}^d$ equals $c\,u\ket0$ for a unitary
    $u$ and a scalar $c\neq0$, and every nonzero product vector is a scalar
    multiple of $(\bigotimes_iu_i)\ket{0\cdots0}$ for one-site unitaries
    $u_i$. A tensor product $\bigotimes_iu_i$ of one-site unitaries is a
    local circuit of depth $2$. Every nonzero vector prepared in depth $T$ is a
    scalar multiple of $U\ket{0\cdots0}$ for a local circuit $U$ of depth
    $T+2$.
\end{lemma}
\begin{proof}
    \leanok
    \uses{lem:qc_circuit_inverse_composition}
    The first layer applies $u_k\otimes u_{k+1}$ on the pairs $\{k,k+1\}$ with
    $k$ odd and $k<N$, and the second applies $u_N$ when $N$ is odd. A nonzero
    product vector $\bigotimes_i\ket{v_i}$ has $v_i\neq0$ for every $i$, and
    $\ket{v_i}=c_iu_i\ket0$ for a unitary $u_i$ whose first column is
    $\ket{v_i}/\lVert v_i\rVert$ and $c_i=\lVert v_i\rVert$. Hence
    $\bigotimes_i\ket{v_i}=(\prod_ic_i)(\bigotimes_iu_i)\ket{0\cdots0}$, and
    Lemma~\ref{lem:qc_circuit_inverse_composition} composes the two
    circuits.
\end{proof}

\begin{lemma}[Operators placed on sets of sites]
    \label{lem:qc_placed_operators}
    \lean{QuantumCircuit.embedOp, QuantumCircuit.embedOp_embedOp,
        QuantumCircuit.embedOp_rectKronecker,
        QuantumCircuit.commute_embedOp_of_disjoint,
        QuantumCircuit.noncommProd_embedOp_apply}
    \leanok
    \uses{def:qc_local_circuit}
    For a map $e$ from the $m$ sites of a chain into the $n$ sites of another
    chain and an operator $X$ on $(\C^d)^{\otimes m}$, the operator $X$
    \emph{placed by} $e$ has the matrix element $X_{x\circ e,\,y\circ e}$ at
    the pair of configurations $(x,y)$ if $x$ and $y$ agree at every site
    outside the range of $e$, and $0$ otherwise; for injective $e$ it is $X$
    acting on the sites $e(1),\dots,e(m)$ and the identity on the other sites.
    Placement is defined in the same way for a map $e\colon J\to I$ of any
    finite sets of sites; for injective $e$, placing a product
    $\bigotimes_{j\in J}m_j$ gives the product with $m_j$ at the site $e(j)$
    and $\mathbb 1$ off $e(J)$. For an injective map $f$, placing by $f$ the
    operator $X$ placed by $e$ places $X$ by $f\circ e$. Operators placed by
    injective maps between chains with disjoint ranges commute. If $e_k$, $k\in K$, are
    injective maps with pairwise disjoint ranges and $X_k$ are operators on
    the corresponding chains, the product of the operators $X_k$ placed by
    $e_k$ has the matrix element $\prod_{k\in K}(X_k)_{x\circ e_k,\,y\circ
            e_k}$ at $(x,y)$ if $x$ and $y$ agree at every site outside all the
    ranges, and $0$ otherwise.
\end{lemma}
\begin{proof}
    \leanok
    A configuration $x$ agrees with $y$ outside the range of $f\circ e$
    exactly when it agrees with $y$ outside the range of $f$ and $x\circ f$
    agrees with $y\circ f$ outside the range of $e$. Placed operators with
    disjoint ranges act on disjoint sets of sites, hence commute. The product
    formula follows by induction on $K$: the sum over intermediate
    configurations in the matrix product has a single nonzero term, the
    configuration agreeing with the output on the range of the new map and
    with the input elsewhere.
\end{proof}

\begin{lemma}[Products of matrices placed on disjoint windows]
    \label{lem:qc_placed_configurations}
    \lean{QuantumCircuit.placeCfg, QuantumCircuit.placeCfg_apply,
        QuantumCircuit.placeCfg_comp, QuantumCircuit.placeCfg_apply_of_notMem,
        QuantumCircuit.eq_placeCfg, QuantumCircuit.placeCfg_injective,
        QuantumCircuit.injective_window, QuantumCircuit.disjoint_range_window,
        QuantumCircuit.list_prod_embedOp_placeCfg, Matrix.mul_apply_extend}
    \leanok
    \uses{lem:qc_placed_operators}
    Let $d\ge1$, and let $W_1,\dots,W_n$ be windows of $m$ sites of a chain,
    that is maps from $\{1,\dots,m\}$ into the sites such that
    $(p,t)\mapsto W_p(t)$ is injective; then every $W_p$ is injective and the
    ranges of distinct windows are disjoint. For configurations $c_p$ of $m$
    sites, the configuration with $c_p$ on the windows carries $c_p(t)$ at
    $W_p(t)$ and $0$ at every other site. It depends injectively on
    $(c_1,\dots,c_n)$, and every configuration with $0$ outside the windows is
    the configuration with its own restrictions on the windows.
    \begin{enumerate}
        \item Let $Y_1,\dots,Y_n$ be matrices on $m$ sites with
              $Y_p\ket{\iota_p(c)}=\sum_eV_p(e,c)\ket{o_p(e)}$ for injective
              maps $o_p$. Then the product of the matrices $Y_p$ placed by the
              windows maps the configuration with $\iota_p(c_p)$ on the windows
              to $\sum_e\prod_pV_p(e_p,c_p)$ times the configuration with
              $o_p(e_p)$ on the windows.
        \item Let $R$ and $R'$ be injective maps into the basis of a space. If
              $M\ket{x_0}=\sum_af(a)\ket{R(a)}$ and
              $G\ket{R(a)}=\sum_bg(a,b)\ket{R'(b)}$ for every $a$, then
              $GM\ket{x_0}=\sum_b\bigl(\sum_ag(a,b)f(a)\bigr)\ket{R'(b)}$.
    \end{enumerate}
\end{lemma}
\begin{proof}
    \leanok
    The statements on configurations are evaluations of the definition. For
    the first item, the placed matrices act on disjoint windows, so by
    Lemma~\ref{lem:qc_placed_operators} the matrix element of their product is
    the product over the windows of the matrix elements of the $Y_p$, and it
    vanishes unless the output configuration is $0$ outside the windows. The
    second item is the matrix product written out.
\end{proof}

\begin{lemma}[Circuits inside sets of sites, in series and in parallel]
    \label{lem:qc_parallel_circuits}
    \lean{QuantumCircuit.BondLayer.IsIn, QuantumCircuit.IsBondCircuitOn,
        QuantumCircuit.IsCircuitOn, QuantumCircuit.IsBondCircuitOn.mul,
        QuantumCircuit.IsBondCircuitOn.finset_noncommProd,
        QuantumCircuit.IsCircuitOn.finset_noncommProd,
        QuantumCircuit.IsPairProduct.isCircuitOn,
        QuantumCircuit.isCircuitOn_list_prod_embedOp}
    \leanok
    \uses{def:qc_local_circuit, lem:qc_circuit_inverse_composition,
        lem:qc_placed_operators}
    A layer acts \emph{inside} a set of sites $R$ if each of its bonds is
    contained in $R$, and $U$ is a \emph{local circuit of depth $T$ inside
        $R$} if it is implemented by $T$ layers acting inside $R$. If $U$ and
    $U'$ are local circuits inside $R$ of depths $T$ and $T'$, then $U'U$ is a
    local circuit inside $R$ of depth $T+T'$. For a bond geometry whose bonds
    are nonempty, such as the ring, if $U_i$, $i\in I$ for a finite set $I$,
    are pairwise commuting local circuits of depth $T$ inside pairwise
    disjoint sets $R_i$, then their product is a local circuit of depth $T$
    inside $\bigcup_{i\in I}R_i$. Let $e$ be an injective map from
    the sites of an open chain of $m$ sites into the ring with
    $e(j+1)=e(j)+1$. If $X$ is a product of at most $K$ unitaries, each acting
    on two neighbouring sites of the open chain, then $X$ placed by $e$ is a
    local circuit of depth $K$ inside the range of $e$. For a list of such
    maps $e_i$ with pairwise disjoint ranges and such products $X_i$, the
    product of the $X_i$ placed by $e_i$ is a local circuit of depth $K$
    inside the union of the ranges.
\end{lemma}
\begin{proof}
    \leanok
    In series the layers of $U$ are followed by those of $U'$
    (Lemma~\ref{lem:qc_circuit_inverse_composition}). In parallel, by
    induction on $I$, the layers of two circuits of depth $T$ inside disjoint
    sets are merged one by one: the union of two layers acting inside disjoint
    sets is a layer, whose operator is the product of the two because the
    nonempty bonds of the two layers are distinct, and the circuit inside one
    set commutes with every layer inside the other. A product of $K$ gates on neighbouring sites of the open
    chain, placed by $e$, is the circuit whose $j$-th layer is the $j$-th gate
    placed by $e$, padded with empty layers; the products on disjoint ranges
    then run in parallel.
\end{proof}

\section{Two-site gates suffice}
\label{sec:qc_two_site_gates}

\begin{theorem}[Unitaries on a bounded number of sites]
    \label{thm:qc_bounded_unitary_two_site_gates}
    \uses{def:qc_local_circuit}
    \lean{QuantumCircuit.exists_isPairProduct,
        QuantumCircuit.isTwoLevelWord_of_det_eq_one,
        QuantumCircuit.exists_isPairProduct_ctrlOp,
        QuantumCircuit.isPairProduct_embedOp_pairSites}
    \leanok
    For $d\ge1$ and $n\ge2$ there is $K$ such that every unitary on
    $(\C^d)^{\otimes n}$ is a product of at most $K$ unitaries, each acting on
    two neighbouring sites of the open chain of $n$ sites. This is the step
    that the unitaries with constant support ``can be further expressed with
    a low-depth circuit of local gates''
    \cite[caption of Fig.~1]{Malz2024Preparation}.
\end{theorem}
\begin{proof}
    \leanok
    Dividing by a root of the determinant of order $d^n$, a global phase and
    hence a single-site gate, reduces to determinant one. Column-by-column Givens elimination writes such
    a unitary as a product of two-level operators, each a real rotation or a
    diagonal phase on the span of two basis vectors. Write $G_{ab}(g)$ for the
    operator acting as the $2\times2$ matrix $g$ on the span of the basis vectors
    $\ket a,\ket b$ and as the identity on their complement, and
    $Q=\bigl(\begin{smallmatrix}0&1\\-1&0\end{smallmatrix}\bigr)$ for the
    quarter turn. For pairwise distinct $a,b,c$,
    \begin{align}
        G_{ab}(g) & =G_{cb}(Q)^\dagger\,G_{ac}(g)\,G_{cb}(Q),
        \label{eq:qc_two_level_quarter_turn}
    \end{align}
    and choosing $c$ to differ from $a$ at one site fewer than $b$ does reduces
    a two-level operator between configurations differing at several sites to
    one between configurations differing at one site fewer. Between configurations
    differing at one site $t$ it is a single-site operation at $t$ controlled by
    the other sites. Every rotation and every phase of determinant one is a group
    commutator $g_1g_2g_1^\dagger g_2^\dagger$. Write $C_S(u)$ for the single-site
    operation $u$ at $t$ controlled by the sites of $S$, with $t\notin S$; for
    unitaries $v,w$ and a site $s\neq t$,
    \begin{align}
        C_{S\cup\{s\}}(vwv^\dagger w^\dagger)
         & =C_S(v)\,C_{\{s\}}(w)\,C_S(v^\dagger)\,C_{\{s\}}(w^\dagger),
        \label{eq:qc_controlled_commutator}
    \end{align}
    so a controlled commutator is a product of two operations with one control
    fewer and two operations on two sites; induction removes the controls. An
    operation on two distant sites is conjugated to neighbouring sites by SWAP
    gates.
\end{proof}

\begin{definition}[Gates on windows of consecutive sites]
    \label{def:qc_window_gate}
    \lean{QuantumCircuit.windowSites, QuantumCircuit.IsWindowGate,
        QuantumCircuit.IsWindowProduct}
    \leanok
    \uses{def:qc_local_circuit}
    On the open chain of $n$ sites, a \emph{gate on a window of $m$ sites} is
    a unitary acting on the sites $p$ with $a\le p<a+m$, for some $a$. Near the
    right end of the chain fewer than $m$ of these sites exist, so such a gate
    acts on at most $m$ consecutive sites. An operator is a \emph{product of at
        most $K$ gates on windows of $m$ sites} if it is a product of at most
    $K$ such gates. These are the gates of the sequential circuit
    of~\cite[paragraph ``The sequential-RG circuit'']{Malz2024Preparation},
    each of which acts on the bond register and one further site.
\end{definition}

\begin{lemma}[Products of gates on windows]
    \label{lem:qc_window_products}
    \lean{QuantumCircuit.isWindowGate_embedOp, QuantumCircuit.isWindowGate_of_le,
        QuantumCircuit.IsWindowProduct.mul, QuantumCircuit.IsWindowProduct.embedOp,
        QuantumCircuit.IsWindowGate.isNeighbourGate,
        QuantumCircuit.IsWindowProduct.isPairProduct_two,
        QuantumCircuit.IsWindowProduct.exists_isPairProduct}
    \leanok
    \uses{def:qc_window_gate, lem:qc_placed_operators,
        thm:qc_bounded_unitary_two_site_gates}
    On the open chain of $n$ sites the following hold.
    \begin{enumerate}
        \item A unitary on $l\le m$ sites placed on consecutive sites is a gate
              on a window of $m$ sites, and if $n\le m$ every unitary on the
              chain is one such gate.
        \item Products of at most $K$ and $K'$ gates on windows of $m$ sites
              multiply to a product of at most $K+K'$ such gates, and placing
              the chain as a block of consecutive sites of a longer chain
              keeps a product of at most $K$ such gates.
        \item If $n\ge2$, a gate on a window of two sites is a gate on two
              neighbouring sites.
        \item For $m\ge2$ there is $C$, depending only on $d$ and $m$, such
              that for every $n\ge2$ a product of at most $K$ gates on windows
              of $m$ sites is a product of at most $KC$ gates on two
              neighbouring sites.
    \end{enumerate}
\end{lemma}
\begin{proof}
    \leanok
    The first two items are the support statements of
    Lemma~\ref{lem:qc_placed_operators}: a placed operator acts on the range
    of the placement, and placing composes. For the third item, the sites of
    the window lie in a pair $\{a',a'+1\}$ with $a'=\min(a,n-2)$. For the
    fourth, enlarge the window of each gate to $\min(m,n)\ge2$ consecutive
    sites inside the chain. The gate is then a unitary on these sites placed
    on the chain, hence by
    Theorem~\ref{thm:qc_bounded_unitary_two_site_gates} a product of at most
    $K_{\min(m,n)}$ gates on neighbouring sites; take $C=\sum_{l\le m}K_l$.
\end{proof}

%% ==================================================
%% Section: Circuits with measurements
%% Sources: arXiv:2307.01696 (Malz, Styliaris, Wei, Cirac), paragraphs
%% "Preparations using measurements" and "Long-range MPS using
%% measurements"; arXiv:2103.13367 (Piroli, Styliaris, Cirac), paragraphs
%% "Quantum circuits and LOCC" and "State transformations with QC and LOCC",
%% Definition "Transformations under QC and LOCC", and Example 1.
%% ==================================================

\section{Circuits with measurements}
\label{sec:qc_measurements}

The model of preparation with measurements is that of \cite{Piroli2021LOCC}:
a local circuit, measurements, and local unitaries depending on the outcomes
of all the measurements, which are communicated to every site at no cost.
Only the layers of the circuit are counted. In this model GHZ-type states are
prepared in constant depth, although their connected correlations do not
decay with distance, while by Lemma~\ref{lem:qc_light_cone_correlations} those
of a state prepared by a local circuit of depth $T$ vanish beyond distance
$2T$. Rounds of measurements are composed, and a vector prepared with
measurements may be followed by a local circuit that does not depend on the
outcomes.

\begin{definition}[Preparation with measurements]
    \label{def:qc_measurement_preparation}
    \lean{QuantumCircuit.outcomeProj, QuantumCircuit.outcomeProj_eq_ctrlProj,
        QuantumCircuit.MeasurementProtocol,
        QuantumCircuit.MeasurementProtocol.preMeasurement,
        QuantumCircuit.MeasurementProtocol.postMeasurement,
        QuantumCircuit.MeasurementProtocol.output,
        QuantumCircuit.MeasurementProtocol.IsPreparationOf,
        QuantumCircuit.IsPreparedWithMeasurementsInDepth}
    \leanok
    \uses{def:qc_local_circuit}
    The local labels are $\{0,\dots,d-1\}$. For a set $S$ of sites of the
    ring of $N$ sites and a string $m\colon S\to\{0,\dots,d-1\}$, let $P_m$
    be the projection onto the span of the basis vectors $\ket{x}$ with
    $x_i=m_i$ for every $i\in S$; these are the outcome projections of a
    measurement of the sites of $S$ in the computational basis, and $P_m$ is
    the projection onto the configurations agreeing on $S$ with any
    configuration extending $m$. A \emph{protocol} consists of a product
    vector $\ket{\pi}$, a local circuit $U$ of depth $T_1$, a set $S$ of
    measured sites, and for every $m$ unitaries $v_{m,i}$ on the sites $i$,
    with product $V_m=\bigotimes_iv_{m,i}$. It \emph{prepares} a vector
    $\ket{\psi}$ if $\ket{\pi}\neq0$ and for every $m$ with
    $P_mU\ket{\pi}\neq0$ there is $c\in\mathbb C$ with
    $V_mP_mU\ket{\pi}=c\ket{\psi}$. A vector is \emph{prepared with
        measurements in depth $T$} if a protocol with $T_1\leq T$ prepares
    it.

    The model is that of \cite{Piroli2021LOCC}, where the class of such
    preparations is called $\mathrm{QCcc}$. The correction $V_m$ may depend on
    the whole string $m$: classical communication is global and free, and the
    single-site corrections are not counted. Ancillas are sites of the ring,
    and the local unitaries between the layers of the circuit, acting on a
    site and its ancillas and free in \cite{Piroli2021LOCC}, are counted here
    as gates. The correction $V_m$ is a local unitary in the sense of
    \cite{Piroli2021LOCC}. So every preparation in this sense is one in the
    sense of \cite{Piroli2021LOCC}, of no larger depth.
\end{definition}

\begin{lemma}[Outcomes of nonzero probability]
    \label{lem:qc_measurement_outcomes}
    \lean{QuantumCircuit.outcomeProj_mulVec_apply,
        QuantumCircuit.sum_outcomeProj, Matrix.finKronecker_mem_unitary,
        Matrix.eq_zero_of_mem_unitary_of_mulVec_eq_zero,
        QuantumCircuit.MeasurementProtocol.sum_postMeasurement,
        QuantumCircuit.MeasurementProtocol.preMeasurement_ne_zero,
        QuantumCircuit.MeasurementProtocol.exists_postMeasurement_ne_zero,
        QuantumCircuit.MeasurementProtocol.output_ne_zero,
        QuantumCircuit.MeasurementProtocol.IsPreparationOf.ne_zero}
    \leanok
    \uses{def:qc_measurement_preparation}
    The outcome projections satisfy $\sum_mP_m=\mathbb 1$, and a product of
    unitaries on the sites is unitary; a unitary matrix maps only the zero
    vector to zero. Hence, for a protocol with
    $\ket{\pi}\neq0$, some outcome has $P_mU\ket{\pi}\neq0$, then
    $V_mP_mU\ket{\pi}\neq0$, and every vector that the protocol prepares is
    nonzero.
\end{lemma}
\begin{proof}
    \leanok
    For each configuration $x$ exactly one string $m$, the restriction of $x$
    to $S$, has $P_m\ket{x}\neq0$. So $\sum_mP_mU\ket{\pi}=U\ket{\pi}$, which
    is nonzero because $U$ is unitary; the unitary $V_m=\bigotimes_iv_{m,i}$,
    with $V_m^\dagger V_m=\bigotimes_iv_{m,i}^\dagger v_{m,i}=\mathbb 1$,
    keeps a nonzero vector nonzero, and $V_mP_mU\ket{\pi}=c\ket{\psi}\neq0$ forces
    $\ket{\psi}\neq0$.
\end{proof}

\begin{lemma}[Protocols without measurements]
    \label{lem:qc_measurement_unitary_case}
    \lean{QuantumCircuit.isPreparedWithMeasurementsInDepth_of_isPreparedInDepth,
        QuantumCircuit.exists_isPreparedInDepth_of_isPreparationOf_of_measured_eq_empty,
        QuantumCircuit.smul_productVector}
    \leanok
    \uses{def:qc_measurement_preparation, lem:qc_measurement_outcomes}
    A nonzero vector prepared in depth $T$ is prepared with measurements in
    depth $T$. Conversely, if a protocol measuring no site, with $T_1\leq T$,
    prepares $\ket{\psi}$, then there are unitaries $u_i$ on the sites $i$
    such that $\bigotimes_iu_i\ket{\psi}$ is prepared in depth $T$.
\end{lemma}
\begin{proof}
    \leanok
    For the first statement take $S=\emptyset$ and $v_{m,i}=\mathbb 1$. For
    the second, $P_m=\mathbb 1$ for the only string $m$, so
    $V_mU\ket{\pi}=c\ket{\psi}$ with $c\neq0$ by
    Lemma~\ref{lem:qc_measurement_outcomes}. Take $u_i=v_{m,i}^\dagger$. Then
    $\bigotimes_iu_i\ket{\psi}=U(c^{-1}\ket{\pi})$, where $c^{-1}\ket{\pi}$
    is the product vector with the vector of one site multiplied by $c^{-1}$,
    and the circuit $U$ of depth at most $T$ is padded with empty layers to
    depth $T$.
\end{proof}

\begin{lemma}[Layers of permutation gates]
    \label{lem:qc_permutation_layers}
    \lean{QuantumCircuit.IsLocalPerm,
        QuantumCircuit.IsLocalPerm.permMatrix_mem_supportedOperators,
        QuantumCircuit.shiftPerm, QuantumCircuit.shiftPerm_apply,
        QuantumCircuit.isLocalPerm_shiftPerm, QuantumCircuit.permLayer,
        QuantumCircuit.exists_permLayer_op_mulVec,
        QuantumCircuit.exists_shiftLayer_op_mulVec,
        Matrix.finKronecker_permMatrix_mulVec}
    \leanok
    \uses{def:qc_local_circuit}
    Let $\sigma$ be a permutation of the configurations of the ring that
    changes only the sites of a set $X$, the new values at these sites
    depending only on the old values there. Its permutation matrix $P_\sigma$,
    $(P_\sigma v)(x)=v(\sigma x)$, is a unitary acting on $X$. Examples are the
    shift $x\mapsto x+f(x)e_t$ at a site $t\in X$, with values in
    $\mathbb Z_d$ and $f$ reading only the sites of $X\setminus\{t\}$. If
    $\sigma_k$, $k\in K$, are such permutations for pairwise disjoint
    pairs $\{k,k+1\}$, the layer of the gates $P_{\sigma_k}$ acts as
    $v\mapsto v\circ\rho$, where $\rho(x)$ agrees with $\sigma_k(x)$ on the
    pair $\{k,k+1\}$ and with $x$ off all the pairs. For shifts
    $\sigma_k\colon x\mapsto x+f_k(x)e_{t_k}$ this reads
    $\rho(x)_{t_k}=x_{t_k}+f_k(x)$, and $\rho(x)_i=x_i$ at every site $i$
    that is none of the $t_k$. Finally, for permutations $\tau_i$ of the
    labels of the sites $i$, the product $\bigotimes_iP_{\tau_i}$ acts as
    $v\mapsto v\circ\tau$, where $\tau(x)_i=\tau_i(x_i)$.
\end{lemma}
\begin{proof}
    \leanok
    Enumerating $X$ writes $P_\sigma$ as $Y\otimes\mathbb 1$ with
    $Y_{u,w}=1$ if $w=\sigma(x)|_X$ for some $x$ with $x|_X=u$, and $0$
    otherwise; permutation matrices are unitary. For the layer, induct over
    the pairs: $P_{\sigma_a}P_\rho=P_{\rho\sigma_a}$, and $\rho\sigma_a$
    agrees with $\sigma_a$ on the pair $\{a,a+1\}$, which the earlier pairs
    do not meet, and with $\sigma_k$ on an earlier pair $\{k,k+1\}$, because
    $\sigma_a(x)$ agrees with $x$ there. The entry of
    $\bigotimes_iP_{\tau_i}$ at $(x,y)$ is $\prod_i[y_i=\tau_i(x_i)]$, which is
    $1$ if $y=\tau(x)$ and $0$ otherwise.
\end{proof}

From here on the local labels are $\mathbb Z_d=\{0,\dots,d-1\}$ with
$d\ge1$.

\begin{definition}[Measurement rounds]
    \label{def:qc_measurement_rounds}
    \lean{QuantumCircuit.MeasurementRound, QuantumCircuit.MeasurementRound.kraus,
        QuantumCircuit.MeasurementRound.depth,
        QuantumCircuit.MeasurementRound.toProtocol,
        QuantumCircuit.MeasurementRound.outputs,
        QuantumCircuit.IsPreparedWithMeasurementRoundsInDepth}
    \leanok
    \uses{def:qc_measurement_preparation}
    A \emph{measurement round} consists of a local circuit $U$, a set $S$ of
    sites measured in the computational basis, and for every outcome
    $m\colon S\to\mathbb Z_d$ unitaries $v_{m,i}$ on the sites $i$, with product
    $V_m=\bigotimes_{i}v_{m,i}$. Its operator for the outcome $m$ is $V_{m}P_{m}U$,
    with $P_m$ the outcome projection of
    Definition~\ref{def:qc_measurement_preparation}, and its depth is the
    number of layers of $U$. With a product vector it is a protocol in the sense
    of that definition. The \emph{outputs} of a sequence of rounds
    $R_1,\dots,R_n$ from a vector $\ket v$ are the vectors
    $K_n\cdots K_1\ket v$, with $K_j$ the operator of $R_j$ for some outcome. A
    vector $\ket\psi$ is \emph{prepared with measurement rounds in depth $T$}
    if there are a product vector $\ket\pi\neq0$ and rounds of total depth at
    most $T$ such that every nonzero output from $\ket\pi$ is a multiple of
    $\ket\psi$.

    Every round is a protocol of $\mathrm{QCcc}_\ell$ of
    \cite{Piroli2021LOCC} of no larger depth, in the sense of
    Definition~\ref{def:qc_measurement_preparation}; that paper performs a
    single round and notes that ``one could also define a more general scheme
    with multiple rounds of LOCC''. A sequence of rounds is a composition of
    such transformations, with the depth counted as the total number of layers.
    The number of rounds is not bounded: the tree circuit of
    Chapter~\ref{ch:log_depth_preparation} with $k$ levels uses $2k$ rounds,
    whereas in the class $\mathrm{QCcc}^{(k)}_\ell$ of the paragraph ``Phases
    of matter'' of \cite{Piroli2021LOCC} the number $k$ of composed
    transformations does not depend on the system size. Several rounds are needed when
    a correction must precede gates that do not commute with it, as for a
    teleported register on which an isometry acts.
\end{definition}

\begin{lemma}[Outputs of rounds]
    \label{lem:qc_measurement_rounds_outputs}
    \lean{QuantumCircuit.MeasurementRound.toProtocol_output,
        QuantumCircuit.MeasurementRound.exists_kraus_mulVec_ne_zero,
        QuantumCircuit.MeasurementRound.outputs_nil,
        QuantumCircuit.MeasurementRound.mem_outputs_cons,
        QuantumCircuit.MeasurementRound.mem_outputs_smul,
        QuantumCircuit.MeasurementRound.mem_outputs_append,
        QuantumCircuit.MeasurementRound.exists_mem_outputs_ne_zero,
        QuantumCircuit.IsPreparedWithMeasurementRoundsInDepth.ne_zero,
        QuantumCircuit.IsPreparedWithMeasurementRoundsInDepth.mono,
        QuantumCircuit.isPreparedWithMeasurementRoundsInDepth_of_isPreparedWithMeasurementsInDepth}
    \leanok
    \uses{def:qc_measurement_rounds, lem:qc_measurement_outcomes}
    The output of the protocol of a round and a product vector $\ket\pi$ for
    the outcome $m$ is $V_{m}P_{m}U\ket\pi$. For a nonzero vector some outcome of a
    round gives a nonzero vector, so a nonzero vector has a nonzero output
    under every sequence of rounds, and a vector prepared with measurement
    rounds is nonzero. The outputs of a multiple $c\ket v$ are the multiples by
    $c$ of the outputs of $\ket v$, and the outputs of a concatenation of two
    sequences are the outputs of the second sequence from the outputs of the
    first. A vector prepared with measurements in depth $T$ is prepared with
    one measurement round in depth $T$.
\end{lemma}
\begin{proof}
    \leanok
    Since $\sum_{m}P_{m}=\mathbb 1$ and $U$ is unitary, some $P_{m}U\ket v$ is
    nonzero, and the unitary $V_m$ keeps it nonzero. The other statements
    follow by induction over the sequence of rounds and from the linearity of
    the operators.
\end{proof}

\begin{definition}[Implementation by rounds]
    \label{def:qc_round_implementation}
    \lean{QuantumCircuit.MeasurementRound.IsImplementationOn,
        QuantumCircuit.MeasurementRound.IsRoundsImplementationOn}
    \leanok
    \uses{def:qc_measurement_rounds}
    Let $E$ be a set of vectors and $W$ a matrix. A round \emph{implements}
    $W$ on $E$ if for every outcome $m$ there is $c\in\mathbb C$ with
    $V_{m}P_{m}U\ket v=cW\ket v$ for every $\ket v\in E$. A sequence of rounds
    \emph{implements} $W$ on $E$ if every output from every $\ket v\in E$ is a
    multiple of $W\ket v$.
\end{definition}

\begin{lemma}[Implementations compose]
    \label{lem:qc_round_implementation_compose}
    \lean{QuantumCircuit.MeasurementRound.IsImplementationOn.exists_mem_outputs,
        QuantumCircuit.MeasurementRound.isRoundsImplementationOn_nil,
        QuantumCircuit.MeasurementRound.IsRoundsImplementationOn.mono,
        QuantumCircuit.MeasurementRound.IsRoundsImplementationOn.append,
        QuantumCircuit.MeasurementRound.IsRoundsImplementationOn.isPreparedWithMeasurementRoundsInDepth}
    \leanok
    \uses{def:qc_round_implementation, lem:qc_measurement_rounds_outputs}
    If a round $R$ implements $W$ on $E$ and $\ket v\in E$, every output of $R$
    followed by rounds $R'$ is a multiple of an output of $R'$ from $W\ket v$.
    The empty sequence implements $\mathbb 1$ on every set, and a sequence
    implementing $W$ on $E$ implements it on every subset of $E$. If a sequence
    $\mathcal R$ implements $W$ on $E$, $W$ maps $E$ into $E'$, and a sequence
    $\mathcal R'$ implements $W'$ on $E'$, then $\mathcal R$ followed by
    $\mathcal R'$ implements $W'W$ on $E$. If $\mathcal R$ implements $W$ on
    $E$ and the product vector $\ket\pi\neq0$ lies in $E$, then $W\ket\pi$ is
    prepared with measurement rounds in the total depth of $\mathcal R$.
\end{lemma}
\begin{proof}
    \leanok
    An output of the concatenation is an output $\ket w$ of $\mathcal R'$ from
    an output $c\,W\ket v$ of $\mathcal R$, hence $\ket w=c\ket{w'}$ with
    $\ket{w'}$ an output of $\mathcal R'$ from $W\ket v\in E'$, and
    $\ket{w'}=c'\,W'W\ket v$. The last statement is the definition of
    preparation with measurement rounds.
\end{proof}

\begin{definition}[Preparation with measurements and a circuit]
    \label{def:qc_measurement_circuit_preparation}
    \lean{QuantumCircuit.IsPreparedWithMeasurementsAndCircuitInDepth}
    \leanok
    \uses{def:qc_measurement_preparation, def:qc_local_circuit}
    A vector $\ket\psi$ is \emph{prepared with measurements and a circuit in
        depth $T$} if $\ket\psi=U\ket\varphi$ for a vector $\ket\varphi$
    prepared with measurements in depth $T_1$ and a local circuit $U$ of depth
    $T_2$, with $T_1+T_2\leq T$. The circuit $U$ does not depend on the
    outcomes of the measurements.
\end{definition}

\begin{lemma}[Layers acting on configurations]
    \label{lem:qc_configuration_layers}
    \lean{QuantumCircuit.apply_eq_ite_of_mulVec_eq_comp, QuantumCircuit.sigmaSite,
        QuantumCircuit.sigmaSite_injective, QuantumCircuit.layerCfg,
        QuantumCircuit.layerCfg_apply, QuantumCircuit.layerCfg_comp,
        QuantumCircuit.layerCfg_apply_of_forall_ne, QuantumCircuit.eq_layerCfg_iff,
        QuantumCircuit.noncommProd_embedOp_mulVec_eq_comp,
        QuantumCircuit.embedOp_mulVec_eq_comp, QuantumCircuit.permOp_mulVec,
        QuantumCircuit.embedOp_mulVec_productVector}
    \leanok
    \uses{def:qc_local_circuit, lem:qc_permutation_layers}
    Let $X_k$ be operators on pairwise disjoint sets $S_k$ of sites, each
    acting by a map $f_k$ of the configurations of $S_k$, that is
    $X_kv=v\circ f_k$ for every vector $v$; then $X_k$ has the entry $1$ at
    $(u,f_k(u))$ and $0$ elsewhere. The product of the $X_k$ acts by the map
    $\rho$ with $\rho(x)|_{S_k}=f_k(x|_{S_k})$ for every $k$ and
    $\rho(x)_i=x_i$ at every site $i$ outside all the $S_k$. A permutation
    $\tau$ of the sites acts as $v\mapsto v\circ(x\mapsto x\circ\tau)$. If $X$
    acts on the sites of $S$ and the product vector $\ket\pi$ is $\ket z$ on
    every site of $S$, then $(X\ket\pi)(y)=X_{y|_S,(z,\dots,z)}\pi'(y)$, where
    $\ket{\pi'}$ is the product vector with the all-ones vector at the sites
    of $S$ and the vectors of $\ket\pi$ elsewhere.
\end{lemma}
\begin{proof}
    \leanok
    Applying $X_k$ to the basis vector $\ket u$ gives the entries. The entry of
    the product at $(x,y)$ is $\prod_k[y|_{S_k}=f_k(x|_{S_k})]$ when $x$ and
    $y$ agree outside the $S_k$, and $0$ otherwise; this is $1$ exactly when
    $y=\rho(x)$. For the product vector, only the configuration $y'$ agreeing
    with $y$ outside $S$ and equal to $(z,\dots,z)$ on $S$ contributes.
\end{proof}

\begin{definition}[GHZ-type states]
    \label{def:qc_ghz_state}
    \lean{QuantumCircuit.ghzState, QuantumCircuit.ghzState_apply}
    \leanok
    Let $b,M\geq0$ and $\alpha\in\mathbb C^b$. The \emph{GHZ-type
        state} on $M$ sites of dimension $b$ is
    $\ket{\chi_M}=\sum_{i=0}^{b-1}\alpha_i\ket{i}^{\otimes M}$. For $M\geq1$
    and the sites numbered $0,\dots,M-1$, its amplitude at a configuration $s$
    is $\alpha_{s_0}$ if all $s_n$ are equal and $0$ otherwise.
\end{definition}

\begin{theorem}[GHZ-type states in constant depth with measurements]
    \label{thm:qc_ghz_measurement}
    \lean{QuantumCircuit.ghzSite, QuantumCircuit.withZeroAncillas,
        QuantumCircuit.ghzInitial, QuantumCircuit.ghzProtocol,
        QuantumCircuit.ghzProtocol_output,
        QuantumCircuit.isPreparedWithMeasurementsInDepth_withZeroAncillas_ghzState}
    \leanok
    \uses{def:qc_measurement_preparation, lem:qc_permutation_layers,
        def:qc_ghz_state}
    Let $M\geq1$, $b\geq1$ and $\alpha\in\mathbb C^b$ nonzero, with the
    labels $\{0,\dots,b-1\}=\mathbb Z_b$. On the ring of $2M$ sites of
    dimension $b$, numbered $0,\dots,2M-1$, place the system qudit
    $n\in\{0,\dots,M-1\}$ at the site $2n$ and its ancilla at the site
    $2n+1$. The GHZ-type state
    $\ket{\chi_M}=\sum_{i\in\mathbb Z_b}\alpha_i\ket{i}^{\otimes M}$ of
    Definition~\ref{def:qc_ghz_state} on the system qudits, with every ancilla in $\ket{0}$, is prepared with
    measurements in depth $2$. In the protocol below every outcome gives
    exactly this state. This is the construction of
    \cite[Example~1]{Piroli2021LOCC}, a preparation by $\mathrm{QCcc}_2$, for
    qudits, and the preparation of $\ket{\chi_M}$ in constant depth stated in
    the
    paragraph ``Long-range MPS using measurements'' of
    \cite{Malz2024Preparation}.
\end{theorem}
\begin{proof}
    \leanok
    The amplitude of $\ket{\chi_M}$ at $s$ is $\alpha_{s_0}$ if all $s_n$ are
    equal and $0$ otherwise. Values are taken in $\mathbb Z_b$. Start from the
    product vector with $\sum_i\alpha_i\ket{i}$ on the system qudit $0$,
    $\sum_i\ket{i}$ on the other system qudits, and $\ket{0}$ on the ancillas.
    A first layer of gates $\ket{k}\ket{a}\mapsto\ket{k}\ket{a+k}$ on the
    system qudit $n$ and its ancilla, and a second layer of gates
    $\ket{a}\ket{k}\mapsto\ket{a-k}\ket{k}$ on the ancilla $n$ and the system
    qudit $n+1$, for $n<M-1$, give
    \begin{align*}
        \sum_{k}\alpha_{k_0}\bigotimes_{n=0}^{M-1}\ket{k_n}
        \bigotimes_{n=0}^{M-2}\ket{k_n-k_{n+1}}\otimes\ket{0},
    \end{align*}
    the last factor being the last ancilla. Measure the ancillas $n<M-1$. The
    outcome $m$ forces $k_{n+1}=k_n-m_n$, so the system qudits are left in
    $\sum_i\alpha_i\bigotimes_n\ket{i-p_n}$ with the prefix sums
    $p_n=m_0+\cdots+m_{n-1}$. The single-site corrections, $X^{p_n}$ on the
    system qudit $n$ and $X^{-m_n}$ on the ancilla $n$, give $\ket{\chi_M}$
    with every ancilla in $\ket{0}$. The two layers are layers of permutation
    gates and the corrections are single-site permutations, and
    Lemma~\ref{lem:qc_permutation_layers} computes their action
    configuration by configuration. The product vector is nonzero because
    $\alpha\neq0$.
\end{proof}

%% ==================================================
%% Section: Teleportation and long-range gates
%% Sources: arXiv:2307.01696 (Malz, Styliaris, Wei, Cirac), eq. (16) and
%% paragraph "Tree-RG circuit with measurements"; arXiv:2103.13367 (Piroli,
%% Styliaris, Cirac), paragraph "State transformations with QC and LOCC" and
%% Example 2.
%% ==================================================

\section{Teleportation and long-range gates}
\label{sec:qc_teleportation}

Registers that are far apart ``can be teleported at neighboring registers with
a constant overhead'' by ``creating nearest-neighbor entangled pairs, then
performing simultaneous measurements, and correcting (without postselection)
based on the measurement outcomes'' \cite{Malz2024Preparation}. This section
proves this for qudits of any dimension: any number of sites is moved at once
along chains of nearest-neighbour hops of any length, and two-site gates
joining the sites $a$ and $a+2L+1$, on pairwise disjoint stretches
$\{a,\dots,a+2L+1\}$, are applied in constant depth for every $L$. Gates at
even separation, or on overlapping stretches, are not covered. Registers of
$s\ge1$ consecutive sites are teleported site by site, and a layer of
unitaries on pairs of distant registers is applied in a depth depending only
on $s$ and on the number of gates on neighbouring sites of each unitary.

Throughout, the sites of the ring are $\{0,\dots,N-1\}$ with addition modulo
$N$, and the local labels are $\mathbb Z_d=\{0,\dots,d-1\}$ with $d\ge1$.

\begin{definition}[Fourier and generalized Pauli matrices]
    \label{def:qc_qudit_pauli}
    \lean{QuantumCircuit.quditRoot, QuantumCircuit.quditFourier,
        QuantumCircuit.quditPauli}
    \leanok
    Let $\zeta=e^{2\pi i/d}$. The Fourier matrix is
    $F_{ab}=d^{-1/2}\zeta^{ab}$ for $a,b\in\mathbb Z_d$, and for
    $z,x\in\mathbb Z_d$ the generalized Pauli matrix $X^xZ^z$ maps $\ket q$ to
    $\zeta^{zq}\ket{q+x}$.
\end{definition}

\begin{lemma}[Fourier and Pauli matrices are unitary]
    \label{lem:qc_qudit_pauli_unitary}
    \lean{QuantumCircuit.isPrimitiveRoot_quditRoot,
        QuantumCircuit.quditFourier_apply,
        QuantumCircuit.quditFourier_mem_unitary,
        QuantumCircuit.quditPauli_apply, QuantumCircuit.quditPauli_mem_unitary}
    \leanok
    \uses{def:qc_qudit_pauli}
    The number $\zeta$ is a primitive $d$-th root of unity, and
    $F$ and $X^xZ^z$ are unitary. The entries are
    $(X^xZ^z)_{pq}=\zeta^{zq}$ if $p-x=q$ and $0$ otherwise.
\end{lemma}
\begin{proof}
    \leanok
    The rows of $F$ are orthonormal because
    $\sum_b\zeta^{(a-c)b}=d\,[a=c]$, the geometric sum of a $d$-th root of
    unity $\zeta^{a-c}$ that differs from $1$ unless $a=c$. The matrix $X^xZ^z$
    is a permutation matrix times a diagonal matrix of unimodular numbers.
\end{proof}

\begin{lemma}[Single-site operators, permutations of sites and sites in the state zero]
    \label{lem:qc_site_permutations}
    \lean{QuantumCircuit.finKronecker_update_one_mulVec_apply,
        QuantumCircuit.finKronecker_update_one_permMatrix_mulVec,
        QuantumCircuit.IsZeroOn, QuantumCircuit.IsZeroOn.mono,
        QuantumCircuit.IsZeroOn.smul,
        QuantumCircuit.IsZeroOn.finKronecker_update_one_mulVec,
        QuantumCircuit.cfgPerm, QuantumCircuit.cfgPerm_apply,
        QuantumCircuit.permMatrix_cfgPerm_mul_rectKronecker,
        QuantumCircuit.permMatrix_cfgPerm_mul_finKronecker,
        QuantumCircuit.IsZeroOn.permMatrix_cfgPerm_mulVec,
        QuantumCircuit.permMatrix_cfgPerm_mul_permMatrix_cfgPerm,
        QuantumCircuit.permMatrix_cfgPerm_one,
        QuantumCircuit.permMatrix_cfgPerm_mul_embedOp_mul,
        QuantumCircuit.IsZeroOn.permMatrix_cfgPerm_mulVec_image,
        QuantumCircuit.IsZeroOn.mulVec_of_mem_supportedOperators,
        QuantumCircuit.list_prod_mem_supportedOperators,
        QuantumCircuit.layerOfListGate, QuantumCircuit.layerOfListGate_key,
        QuantumCircuit.layerOfList, QuantumCircuit.layerOfList_op}
    \leanok
    \uses{def:qc_local_circuit}
    A vector $\ket v$ has $\ket0$ at the sites of a set $Z$ if $v(x)=0$ for
    every configuration $x$ with $x_i\neq0$ for some $i\in Z$. A permutation
    $\pi$ of the sites acts on vectors by $(P_\pi v)(x)=v(x\circ\pi)$. These
    two notions, the second and third items below, and the statement on
    scalars in the first item hold for any finite set of sites; the other
    statements are for the ring. Then:
    \begin{enumerate}
        \item the operator $u$ at the site $t$ maps $\ket v$ to
              $x\mapsto\sum_{j}u_{x_t,j}\,v(x[t\mapsto j])$, and keeps $\ket0$ at
              the sites of $Z$ if $t\notin Z$; so does multiplication by a
              scalar;
        \item $P_\pi P_\sigma=P_{\pi\sigma}$, $P_\pi(\bigotimes_{i}g_{i})
                  =(\bigotimes_{i}g_{\pi^{-1}(i)})P_\pi$, and conjugating an operator
              $X$ placed at the sites $e_1,\dots,e_m$ by $P_\sigma$ places $X$
              at the sites $\sigma(e_1),\dots,\sigma(e_m)$:
              $P_\sigma XP_{\sigma^{-1}}$ is $X$ at $\sigma\circ e$;
        \item if $\ket v$ has $\ket0$ at the sites of $Z$, then $P_\pi\ket v$
              has $\ket0$ at the sites of $\pi(Z)$, and an operator acting on
              sites outside $Z$, in particular a product of such operators,
              keeps $\ket0$ at the sites of $Z$;
        \item gates on pairwise disjoint neighbouring pairs, given by a list,
              form a layer whose operator is the product of the gates along the
              list.
    \end{enumerate}
\end{lemma}
\begin{proof}
    \leanok
    These are entrywise computations. For the third item, $\ket v$ has $\ket0$
    on $Z$ exactly when the diagonal projection onto the configurations
    vanishing on $Z$ fixes $\ket v$, and that projection commutes with an
    operator acting outside $Z$.
\end{proof}

\begin{lemma}[One teleportation hop]
    \label{lem:qc_teleport_hop}
    \lean{QuantumCircuit.TeleportHop, QuantumCircuit.TeleportHop.gate₁,
        QuantumCircuit.TeleportHop.gate₂,
        QuantumCircuit.TeleportHop.gate₁_mem_unitary,
        QuantumCircuit.TeleportHop.gate₂_mem_unitary,
        QuantumCircuit.TeleportHop.gate₁_mem_supportedOperators,
        QuantumCircuit.TeleportHop.gate₂_mem_supportedOperators,
        QuantumCircuit.TeleportHop.pre, QuantumCircuit.TeleportHop.frame,
        QuantumCircuit.TeleportHop.frame_mem_unitary,
        QuantumCircuit.TeleportHop.frame_of_ne,
        QuantumCircuit.TeleportHop.swapPerm,
        QuantumCircuit.TeleportHop.pre_mulVec,
        QuantumCircuit.TeleportHop.pre_mem_supportedOperators}
    \leanok
    \uses{def:qc_qudit_pauli, lem:qc_qudit_pauli_unitary,
        lem:qc_site_permutations, lem:qc_permutation_layers}
    A \emph{hop} consists of three distinct sites $c,e,f$ such that $\{e,f\}$
    and $\{c,e\}$ are neighbouring pairs. Its first gate, on $\{e,f\}$,
    applies $F$ at $e$ and then $\ket a_e\ket b_f\mapsto\ket a_e\ket{b+a}_f$;
    its second gate, on $\{c,e\}$, applies
    $\ket a_c\ket b_e\mapsto\ket a_c\ket{b-a}_e$ and then $F$ at $c$. Both are
    unitary and act on their pairs. Let $\ket v$ have $\ket0$ at $e$ and $f$,
    and let $z_c,z_e\in\mathbb Z_d$. Applying the two gates and then the
    projection onto $\ket{z_c}$ at $c$ and $\ket{z_e}$ at $e$ gives
    \begin{align*}
        \frac1d\,\bigl(s_{z_c}\bigr)_c\bigl(s_{z_e}\bigr)_e
        \bigl(X^{z_e}Z^{z_c}\bigr)_{f}P_{(c\,f)}\ket v,
    \end{align*}
    where $s_z\ket q=\ket{q+z}$ and $(c\,f)$ is the transposition of the
    sites $c$ and $f$. So the content of $c$ is found at $f$, up to the
    unitary $X^{z_e}Z^{z_c}$, and $c$ and $e$ carry the outcomes. The operator
    acts on the sites $c,e,f$.
\end{lemma}
\begin{proof}
    \leanok
    On a vector with $\ket0$ at $e$ and $f$, the Fourier matrix at $e$ gives
    the amplitude $d^{-1/2}$ to every value $x_e$, and the shift copies it to
    $f$: this prepares $d^{-1/2}\sum_j\ket{j}_e\ket j_f$. The second gate and
    the projections compute, configuration by configuration, the amplitude
    $d^{-1/2}F_{z_c,\,x_f-z_e}$ for the content $x_f-z_e$ moved from $c$ to
    $f$, which is $d^{-1}\zeta^{z_c(x_f-z_e)}$, the amplitude of the right-hand
    side. Configurations with $x_c\neq z_c$ or $x_e\neq z_e$ vanish on both
    sides.
\end{proof}

\begin{theorem}[Teleportation along chains of hops in one round]
    \label{thm:qc_teleportation_round}
    \lean{QuantumCircuit.TeleportHop.sites,
        QuantumCircuit.TeleportHop.allSites,
        QuantumCircuit.TeleportHop.pairSites,
        QuantumCircuit.TeleportHop.measuredSites,
        QuantumCircuit.TeleportHop.Valid,
        QuantumCircuit.TeleportHop.mem_allSites,
        QuantumCircuit.TeleportHop.coe_measuredSites_subset_allSites,
        QuantumCircuit.TeleportHop.disjoint_measuredSites,
        QuantumCircuit.TeleportHop.pairwise_disjoint_one,
        QuantumCircuit.TeleportHop.pairwise_disjoint_two,
        QuantumCircuit.TeleportHop.layerOne, QuantumCircuit.TeleportHop.layerTwo,
        QuantumCircuit.TeleportHop.layerOne_op,
        QuantumCircuit.TeleportHop.layerTwo_op,
        QuantumCircuit.TeleportHop.sitePerm,
        QuantumCircuit.TeleportHop.chainPerm,
        QuantumCircuit.TeleportHop.chainPerm_nil,
        QuantumCircuit.TeleportHop.chainPerm_cons,
        QuantumCircuit.TeleportHop.chainFrame,
        QuantumCircuit.TeleportHop.chainPre,
        QuantumCircuit.TeleportHop.chainFrame_mem_unitary,
        QuantumCircuit.TeleportHop.chainFrame_of_notMem,
        QuantumCircuit.TeleportHop.prod_gate₂_mem_supportedOperators,
        QuantumCircuit.TeleportHop.chainPre_cons,
        QuantumCircuit.IsZeroOn.chainPerm_mulVec,
        QuantumCircuit.TeleportHop.chainPre_mulVec,
        QuantumCircuit.TeleportHop.extendOutcome,
        QuantumCircuit.TeleportHop.round, QuantumCircuit.TeleportHop.depth_round,
        QuantumCircuit.TeleportHop.isImplementationOn_round}
    \leanok
    \uses{lem:qc_teleport_hop, def:qc_round_implementation,
        lem:qc_site_permutations}
    A list $h_1,\dots,h_H$ of hops is \emph{valid} if every hop $h_j$ meets the
    hops $h_{j+1},\dots,h_H$ after it in the list only through its site $c$:
    its sites $e$ and $f$ are not sites of $h_{j+1},\dots,h_H$, and its site
    $c$ is not a site $c$ or $e$ of one of them. Its site $c$ may be the site
    $f$ of one of them, so the hops form chains
    $c_0\to f_0=c_1\to f_1=\cdots$, the hops after $h_j$ in the list acting
    before it. For a valid list and a local circuit $U$,
    the \emph{teleportation round} applies $U$, the first gates of all hops in
    one layer, the second gates of all hops in one layer, measures the sites
    $c$ and $e$ of all hops, and corrects by the inverse of the product of the
    single-site unitaries that the hops leave. Its depth is the depth of $U$
    plus $2$. It implements $P_{h_1}\cdots P_{h_H}U$, with $P_h$ the
    transposition of the sites $c$ and $f$ of $h$, on the vectors $\ket v$ such
    that $U\ket v$ has $\ket0$ at the sites $e$ and $f$ of every hop, with the
    scalar $d^{-H}$ for every outcome.
\end{theorem}
\begin{proof}
    \leanok
    The first gates act on pairwise disjoint pairs $\{e,f\}$, and the second
    gates on pairwise disjoint pairs $\{c,e\}$, so they form two layers. By
    induction over the list, the projections and the two layers factor into
    the hops one after the other: the first gate of the head $h$ of the list
    commutes with the second gates of the other hops because its sites $e,f$
    are not theirs, and the projections at the sites of the other hops commute
    with the operator of $h$. Lemma~\ref{lem:qc_teleport_hop} for $h$ applies
    because the other hops leave $\ket0$ at $e$ and $f$. A single-site unitary
    at the site $c$ of $h$, left by another hop whose site $f$ is $c$,
    reappears at the site $f$ of $h$ after the transposition; so the hops leave
    a product of single-site unitaries, depending on the outcomes, after
    $d^{-H}P_{h_1}\cdots P_{h_H}$. The correction is its inverse.
\end{proof}

\begin{lemma}[Chains between distant sites]
    \label{lem:qc_teleportation_chains}
    \lean{QuantumCircuit.add_natCast_injective, QuantumCircuit.add_natCast_succ,
        QuantumCircuit.add_natCast_ne,
        QuantumCircuit.TeleportHop.exists_eq_smul_of_mem_outputs_conj,
        QuantumCircuit.TeleportHop.chainHop,
        QuantumCircuit.TeleportHop.chainHopBack,
        QuantumCircuit.TeleportHop.forwardChain,
        QuantumCircuit.TeleportHop.backwardChain,
        QuantumCircuit.TeleportHop.After,
        QuantumCircuit.TeleportHop.valid_iff_pairwise,
        QuantumCircuit.TeleportHop.sitePerm_append,
        QuantumCircuit.TeleportHop.pairSites_append,
        QuantumCircuit.TeleportHop.mem_forwardChain,
        QuantumCircuit.TeleportHop.mem_backwardChain,
        QuantumCircuit.TeleportHop.valid_forwardChain,
        QuantumCircuit.TeleportHop.valid_backwardChain,
        QuantumCircuit.TeleportHop.sitePerm_forwardChain,
        QuantumCircuit.TeleportHop.sitePerm_backwardChain,
        QuantumCircuit.TeleportHop.pairSites_forwardChain_subset,
        QuantumCircuit.TeleportHop.pairSites_backwardChain_subset,
        QuantumCircuit.TeleportHop.allSites_forwardChain_subset,
        QuantumCircuit.TeleportHop.allSites_backwardChain_subset,
        QuantumCircuit.TeleportHop.isZeroOn_chainPerm_forwardChain}
    \leanok
    \uses{thm:qc_teleportation_round, lem:qc_round_implementation_compose}
    \begin{enumerate}
        \item Let $\mathcal T$ and $\mathcal B$ be valid lists of hops and $G$ a
              layer of gates, and let $P_{\mathcal T}$, $P_{\mathcal B}$ be the
              products of the transpositions of their hops. If $\ket v$ has
              $\ket0$ at the sites $e,f$ of $\mathcal T$ and
              $GP_{\mathcal T}\ket v$ has $\ket0$ at the sites $e,f$ of
              $\mathcal B$, then every output of the teleportation round of
              $\mathcal T$, followed by the teleportation round of $\mathcal B$
              after $G$, is a multiple of $P_{\mathcal B}GP_{\mathcal T}\ket v$.
              The two rounds have depth $2$ and $3$.
        \item Let $2L<N$ and let $a$ be a site. The forward chain of the hops
              $(a+2j,\,a+2j+1,\,a+2j+2)$, $j<L$, and the backward chain of the
              hops $(a+2j+2,\,a+2j+1,\,a+2j)$, $j<L$, are valid. The product of
              the transpositions of the forward chain maps $a$ to $a+2L$ and
              fixes the sites $a+j$ with $2L<j<N$; that of the backward chain is
              its inverse. The sites $e,f$ of the forward chain lie among
              $a+1,\dots,a+2L$, those of the backward chain lie among
              $a,\dots,a+2L-1$, and all sites of both chains among
              $a,\dots,a+2L$. If $\ket v$ has $\ket0$ at the sites
              $a+1,\dots,a+2L$, then after the forward chain the sites
              $a,\dots,a+2L-1$ carry $\ket0$.
    \end{enumerate}
\end{lemma}
\begin{proof}
    \leanok
    The first item composes two applications of
    Theorem~\ref{thm:qc_teleportation_round}. For the second, the sites
    $a+j$ with $0\le j<N$ are distinct, the transposition of the hop $j$
    exchanges $a+2j$ and $a+2j+2$, and the hops of the forward chain, applied
    in the order $j=0,1,\dots$, move the content of $a$ step by step to
    $a+2L$; each hop moves the $\ket0$ from $a+2j+2$ to $a+2j$, while
    $a+2j+1$ keeps its $\ket0$.
\end{proof}

\begin{lemma}[Parallel chains]
    \label{lem:qc_parallel_chains}
    \lean{QuantumCircuit.TeleportHop.sitePerm_apply_of_notMem,
        QuantumCircuit.TeleportHop.mem_allSites_flatMap,
        QuantumCircuit.TeleportHop.mem_pairSites_flatMap,
        QuantumCircuit.TeleportHop.chainPerm_append,
        QuantumCircuit.TeleportHop.after_of_disjoint,
        QuantumCircuit.TeleportHop.sites_subset_allSites,
        QuantumCircuit.TeleportHop.valid_flatMap,
        QuantumCircuit.TeleportHop.sitePerm_flatMap_apply,
        QuantumCircuit.TeleportHop.sitePerm_flatMap_apply_of_notMem}
    \leanok
    \uses{lem:qc_teleportation_chains}
    The product of the transpositions of a list of hops fixes every site that
    is not a site of the list. Let valid lists of hops $\mathcal L_g$ have
    their sites in sets $S_g$ that are pairwise disjoint. Then the
    concatenation of the lists is valid, and the product of its transpositions
    agrees on $S_g$ with that of $\mathcal L_g$ and fixes every site outside
    all the $S_g$.
\end{lemma}
\begin{proof}
    \leanok
    Two hops with disjoint sites satisfy the condition of validity in either
    order. The product of the transpositions of $\mathcal L_g$ fixes the sites
    outside $S_g$, hence maps $S_g$ into itself, and the products of the other
    lists fix $S_g$ pointwise.
\end{proof}

\begin{theorem}[A layer of distant gates in depth five with measurements]
    \label{thm:qc_long_range_layer}
    \lean{QuantumCircuit.LongRangeGate, QuantumCircuit.LongRangeGate.far,
        QuantumCircuit.LongRangeGate.near, QuantumCircuit.LongRangeGate.span,
        QuantumCircuit.LongRangeGate.interior,
        QuantumCircuit.LongRangeGate.cleared, QuantumCircuit.LongRangeGate.op,
        QuantumCircuit.LongRangeGate.there, QuantumCircuit.LongRangeGate.back,
        QuantumCircuit.LongRangeGate.near_ne_far,
        QuantumCircuit.LongRangeGate.localGate,
        QuantumCircuit.LongRangeGate.localGate_mem_unitary,
        QuantumCircuit.LongRangeGate.localGate_mem_supportedOperators,
        QuantumCircuit.LongRangeGate.a_ne_far,
        QuantumCircuit.LongRangeGate.op_mem_unitary,
        QuantumCircuit.LongRangeGate.op_mem_supportedOperators,
        QuantumCircuit.LongRangeGate.far_mem_span,
        QuantumCircuit.LongRangeGate.bond_near_subset_span,
        QuantumCircuit.LongRangeGate.interior_subset_span,
        QuantumCircuit.LongRangeGate.cleared_subset_span,
        QuantumCircuit.LongRangeGate.allSites_there_subset_span,
        QuantumCircuit.LongRangeGate.allSites_back_subset_span,
        QuantumCircuit.LongRangeGate.disjoint_cleared_bond_near,
        QuantumCircuit.LongRangeGate.sitePerm_back_near,
        QuantumCircuit.LongRangeGate.sitePerm_back_far,
        QuantumCircuit.LongRangeGate.sitePerm_back_sitePerm_there,
        QuantumCircuit.LongRangeGate.pairwise_disjoint_bond,
        QuantumCircuit.LongRangeGate.localLayer,
        QuantumCircuit.LongRangeGate.localLayer_op,
        QuantumCircuit.LongRangeGate.valid_there,
        QuantumCircuit.LongRangeGate.valid_back,
        QuantumCircuit.LongRangeGate.isZeroOn_chainPerm_there,
        QuantumCircuit.LongRangeGate.isZeroOn_localLayer,
        QuantumCircuit.LongRangeGate.sitePerm_there_eq_symm,
        QuantumCircuit.LongRangeGate.chainPerm_back_mul_localLayer_mul_chainPerm_there,
        QuantumCircuit.LongRangeGate.rounds,
        QuantumCircuit.LongRangeGate.sum_depth_rounds,
        QuantumCircuit.LongRangeGate.isRoundsImplementationOn_rounds}
    \leanok
    \uses{lem:qc_teleportation_chains, lem:qc_parallel_chains}
    A \emph{distant gate} is a two-site unitary $u$ on the sites $a$ and
    $a+2L+1$, where $2L+1<N$; its stretch is $a,a+1,\dots,a+2L+1$ and its
    interior the $2L$ sites $a+1,\dots,a+2L$. Let $g_1,\dots,g_n$ be distant
    gates with pairwise disjoint stretches. Two measurement rounds of total
    depth $5$ implement the product of the gates $g_1,\dots,g_n$ on the
    vectors with $\ket0$ at the interiors of all the gates: the first
    teleports the content of every site $a$ along its forward chain to
    $a+2L$; the second applies every $u$ to the neighbouring pair
    $\{a+2L,a+2L+1\}$ and teleports the content of every $a+2L$ back to $a$
    along its backward chain. The depth depends neither on the distances nor
    on the number of gates, and no outcome is post-selected.
\end{theorem}
\begin{proof}
    \leanok
    By Lemma~\ref{lem:qc_parallel_chains} the forward chains of all the gates
    form a valid list, and so do the backward chains, with products of
    transpositions $P$ and $P^{-1}$: on the stretch of $g$ both agree with the
    chains of $g$, which are inverse to each other. The forward chains need
    $\ket0$ at the interiors; afterwards the sites $a,\dots,a+2L-1$ of every
    gate carry $\ket0$, by Lemma~\ref{lem:qc_teleportation_chains} for its
    own chain and since the other chains act on disjoint sites. The gates on
    the pairs $\{a+2L,a+2L+1\}$ keep these $\ket0$, which the backward chains
    need. By Lemma~\ref{lem:qc_teleportation_chains} every output is a
    multiple of $P^{-1}GP\ket v$, with $G$ the product of the gates on the
    neighbouring pairs, and $P^{-1}GP$ is the product of the conjugated gates.
    The conjugate of $u$ at $a+2L,a+2L+1$ is $u$ at $P^{-1}(a+2L)=a$ and
    $P^{-1}(a+2L+1)=a+2L+1$ by Lemma~\ref{lem:qc_site_permutations}.
\end{proof}

\begin{definition}[Register gates]
    \label{def:qc_register_gate}
    \lean{QuantumCircuit.RegisterGate, QuantumCircuit.RegisterGate.offset,
        QuantumCircuit.RegisterGate.sites,
        QuantumCircuit.RegisterGate.localSites,
        QuantumCircuit.RegisterGate.span, QuantumCircuit.RegisterGate.zone,
        QuantumCircuit.RegisterGate.interior, QuantumCircuit.RegisterGate.op,
        QuantumCircuit.RegisterGate.localOp,
        QuantumCircuit.RegisterGate.offset_lt,
        QuantumCircuit.RegisterGate.sites_injective,
        QuantumCircuit.RegisterGate.localSites_injective,
        QuantumCircuit.RegisterGate.localSites_succ,
        QuantumCircuit.RegisterGate.sites_mem_span,
        QuantumCircuit.RegisterGate.localSites_mem_span,
        QuantumCircuit.RegisterGate.zone_subset_span,
        QuantumCircuit.RegisterGate.op_mem_unitary,
        QuantumCircuit.RegisterGate.op_mem_supportedOperators,
        QuantumCircuit.RegisterGate.add_eq_add_iff,
        QuantumCircuit.RegisterGate.two_mul_L_lt}
    \leanok
    \uses{def:qc_local_circuit, lem:qc_placed_operators}
    A \emph{register gate} is a matrix $X$ on $2s$ sites together with a site
    $a$ and $L\ge0$ with $2L+2s\le N$. It acts on the two registers
    $a,\dots,a+s-1$ and $a+2L+s,\dots,a+2L+2s-1$, the first register on the
    first $s$ sites of $X$. Its stretch is $a,\dots,a+2L+2s-1$, its interior the
    $2L$ sites $a+s,\dots,a+s+2L-1$ strictly between the registers, and, for
    $0\le t\le s$, its zone $t$ the $2L$ sites $a+t,\dots,a+t+2L-1$. The gate
    acts on the ring as $X$ on its two registers, which is unitary when $X$ is.
\end{definition}

\begin{lemma}[Teleporting registers site by site]
    \label{lem:qc_register_chains}
    \lean{QuantumCircuit.IsZeroOn.permMatrix_cfgPerm_mulVec_of_mapsTo,
        Equiv.Perm.apply_mem_of_forall_notMem,
        QuantumCircuit.RegisterGate.there, QuantumCircuit.RegisterGate.back,
        QuantumCircuit.RegisterGate.allSites_there_subset,
        QuantumCircuit.RegisterGate.allSites_back_subset,
        QuantumCircuit.RegisterGate.allSites_there_subset_span,
        QuantumCircuit.RegisterGate.allSites_back_subset_span,
        QuantumCircuit.RegisterGate.pairSites_there_subset,
        QuantumCircuit.RegisterGate.pairSites_back_subset,
        QuantumCircuit.RegisterGate.sitePerm_there_apply_of_lt,
        QuantumCircuit.RegisterGate.sitePerm_there_apply_self,
        QuantumCircuit.RegisterGate.sitePerm_back_eq_inv,
        QuantumCircuit.RegisterGate.sitePerm_back_apply_of_lt,
        QuantumCircuit.RegisterGate.sitePerm_back_apply_self,
        QuantumCircuit.RegisterGate.sitePerm_there_notMem_zone,
        QuantumCircuit.RegisterGate.sitePerm_back_notMem_zone,
        QuantumCircuit.RegisterGate.thereAll,
        QuantumCircuit.RegisterGate.backAll,
        QuantumCircuit.RegisterGate.zoneAll,
        QuantumCircuit.RegisterGate.valid_thereAll,
        QuantumCircuit.RegisterGate.valid_backAll,
        QuantumCircuit.RegisterGate.sitePerm_thereAll_apply,
        QuantumCircuit.RegisterGate.sitePerm_backAll_apply,
        QuantumCircuit.RegisterGate.isZeroOn_chainPerm_thereAll,
        QuantumCircuit.RegisterGate.isZeroOn_chainPerm_backAll,
        QuantumCircuit.RegisterGate.pairSites_thereAll_subset,
        QuantumCircuit.RegisterGate.pairSites_backAll_subset,
        QuantumCircuit.RegisterGate.sitePerm_backAll_eq_inv,
        QuantumCircuit.RegisterGate.therePerm,
        QuantumCircuit.RegisterGate.therePerm_apply,
        QuantumCircuit.RegisterGate.therePerm_comp_sites,
        QuantumCircuit.TeleportHop.valid_nil,
        QuantumCircuit.RegisterGate.thereRounds,
        QuantumCircuit.RegisterGate.backRounds,
        QuantumCircuit.RegisterGate.sum_depth_thereRounds,
        QuantumCircuit.RegisterGate.sum_depth_backRounds,
        QuantumCircuit.MeasurementRound.IsImplementationOn.isRoundsImplementationOn,
        QuantumCircuit.RegisterGate.isRoundsImplementationOn_thereRounds,
        QuantumCircuit.RegisterGate.isRoundsImplementationOn_backRounds}
    \leanok
    \uses{def:qc_register_gate, thm:qc_teleportation_round,
        lem:qc_teleportation_chains, lem:qc_parallel_chains,
        lem:qc_site_permutations, lem:qc_round_implementation_compose}
    Let $g_1,\dots,g_n$ be register gates with pairwise disjoint stretches. For
    $0\le t<s$, the round $t$ teleports, for every gate, the content of the site
    $a+t$ to $a+2L+t$ along a chain of $L$ hops, in one round of depth $2$. The
    rounds $t=s-1,\dots,0$, in this order, implement the product $P$ of the
    transpositions of all the chains on the vectors with $\ket0$ at the
    interiors of all the gates, and their outputs have $\ket0$ at the zones $0$.
    The rounds teleporting back, $t=0,\dots,s-1$, in this order, implement
    $P^{-1}$ on the vectors with $\ket0$ at the zones $0$. The rounds of each
    direction have total depth $2s$.
\end{lemma}
\begin{proof}
    \leanok
    The chain of the site $a+t$ meets the sites $a+t,\dots,a+t+2L$, and before
    it the sites $a+t+1,\dots,a+t+2L$ form the zone $t+1$. Starting from $\ket0$
    at the zone $s$, the interior, the round $t$ moves the content of $a+t$ to
    $a+t+2L$ and leaves $\ket0$ at the zone $t$, since the permutation of the
    chain maps the zone $t+1$ together with $a+t$ onto the zone $t$ together
    with $a+t+2L$. The chains of different gates act on disjoint stretches, so
    the lists of all the gates are valid
    (Lemma~\ref{lem:qc_parallel_chains}), and by
    Theorem~\ref{thm:qc_teleportation_round} every round implements the
    permutation of its chains. Composition
    (Lemma~\ref{lem:qc_round_implementation_compose}) gives the claim; the
    backward rounds run the forward ones in reverse.
\end{proof}

\begin{theorem}[A layer of register gates in constant depth with measurements]
    \label{thm:qc_register_layer}
    \lean{List.mul_prod_mul_of_mul_eq_one,
        QuantumCircuit.RegisterGate.rounds,
        QuantumCircuit.RegisterGate.sum_depth_rounds,
        QuantumCircuit.RegisterGate.localProd,
        QuantumCircuit.RegisterGate.isZeroOn_localProd_mulVec,
        QuantumCircuit.RegisterGate.permMatrix_mul_localProd_mul_permMatrix,
        QuantumCircuit.RegisterGate.isZeroOn_zoneAll_iff,
        QuantumCircuit.RegisterGate.isRoundsImplementationOn_rounds,
        QuantumCircuit.RegisterGate.isCircuitOn_localProd,
        QuantumCircuit.RegisterGate.exists_rounds}
    \leanok
    \uses{lem:qc_register_chains, lem:qc_round_implementation_compose}
    Let $g_1,\dots,g_n$ be register gates with pairwise disjoint stretches, each
    $X$ a product of at most $K$ unitaries on neighbouring sites. Then
    measurement rounds of total depth $4s+K+2$ implement the product of the
    gates on the vectors with $\ket0$ at the interiors of all the gates: the
    first registers are teleported next to the second ones as in
    Lemma~\ref{lem:qc_register_chains}, every $X$ is applied to the $2s$
    consecutive sites $a+2L,\dots,a+2L+2s-1$ by a local circuit of depth $K$,
    and the first registers are teleported back. The depth depends neither on
    the distances nor on the number of gates, and no outcome is post-selected.
\end{theorem}
\begin{proof}
    \leanok
    The unitaries on the consecutive sites act inside the stretches, so they
    keep $\ket0$ at the zones $0$, and their product $G$ is a local circuit of
    depth $K$, applied in one round without hops of depth $K+2$. By
    Lemma~\ref{lem:qc_register_chains} every output is a multiple of
    $P^{-1}GP\ket v$. Conjugation by $P$ maps the consecutive sites of every
    gate back to its two registers, so $P^{-1}GP$ is the product of the gates.
\end{proof}

%% ==================================================
%% Section: Local channel circuits and conversions
%% Sources: arXiv:2307.01696 (Malz, Styliaris, Wei, Cirac), proof of
%% Theorem 1 in the Supplemental Material; arXiv:2103.13367 (Piroli,
%% Styliaris, Cirac), paragraphs "Quantum circuits and LOCC" and "State
%% transformations with QC and LOCC".
%% ==================================================

\section{Local channel circuits and conversions}
\label{sec:qc_channels}

Lemmas~\ref{lem:qc_light_cone} and~\ref{lem:qc_light_cone_correlations}
extend from unitary gates to channels, the operations
of a local circuit that may add product ancillas and discard them
\cite{Piroli2021LOCC}.

\begin{definition}[Local channel circuit of depth $T$]
    \label{def:qc_local_channel_circuit}
    \uses{def:qc_local_circuit}
    \lean{QuantumCircuit.ChannelLayer, QuantumCircuit.BondChannelLayer,
        QuantumCircuit.BondChannelLayer.map,
        QuantumCircuit.BondChannelLayer.dual, QuantumCircuit.channelCircuitMap,
        QuantumCircuit.channelCircuitDual,
        QuantumCircuit.IsChannelPreparedInDepth,
        QuantumCircuit.IsBondChannelPreparedInDepth}
    \leanok
    Fix a bond geometry, for instance the pairs $\{k,k+1\}$ of the ring. A
    \emph{channel on the bond} $B$ is a map
    $\rho\mapsto\sum_jK_j\rho K_j^\dagger$ on density matrices of all the
    sites, where each $K_j$ is an operator on all the sites acting trivially
    off $B$ and
    $\sum_jK_j^\dagger K_j=\mathbb 1$. A \emph{channel layer} is a composition
    of channels on disjoint bonds, and a
    \emph{local channel circuit of depth $T$} is a composition
    $\Phi=\Phi_T\circ\dots\circ\Phi_1$ of $T$ channel layers.
    Its dual $\Phi^\dagger=\Phi_1^\dagger\circ\dots\circ\Phi_T^\dagger$
    replaces each channel on a bond by $A\mapsto\sum_jK_j^\dagger AK_j$. A
    density matrix is \emph{prepared in depth $T$ by local channels} if it
    equals $\Phi(\bigotimes_i\sigma_i)$ for a local channel circuit $\Phi$ of
    depth $T$ and density matrices $\sigma_i$ on single sites.
\end{definition}

\begin{lemma}[Duality and complete positivity of local channel circuits]
    \label{lem:qc_local_channel_circuit_properties}
    \uses{def:qc_local_channel_circuit, def:qc_local_circuit}
    \lean{QuantumCircuit.channelCircuitMap_isKrausCPTP,
        QuantumCircuit.trace_channelCircuitMap_mul,
        QuantumCircuit.channelCircuitMap_map_toChannelLayer}
    \leanok
    Let $\Phi$ be a local channel circuit. Then
    \begin{align}
        \Tr(\Phi(\rho)A) & =\Tr(\rho\,\Phi^\dagger(A))
        \label{eq:qc_channel_circuit_duality}
    \end{align}
    for all operators $\rho$ and $A$ on the sites, and $\Phi$ is trace
    preserving and completely positive. A local circuit $U$ of depth $T$ is
    the local channel circuit $\rho\mapsto U\rho U^\dagger$ of depth $T$ with
    one operator per channel.
\end{lemma}
\begin{proof}
    \leanok
    For one channel on a bond,
    $\Tr\bigl(\sum_jK_j\rho K_j^\dagger A\bigr)
        =\Tr\bigl(\rho\sum_jK_j^\dagger AK_j\bigr)$ by cyclicity of the
    trace, and $\sum_jK_j^\dagger K_j=\mathbb 1$ makes it trace preserving
    and completely positive. Both properties pass to compositions, hence to
    layers and circuits. A gate $G$ on a bond is the channel with the single
    operator $G$.
\end{proof}

\begin{lemma}[Light cone of a local channel circuit]
    \label{lem:qc_channel_light_cone}
    \uses{def:qc_local_channel_circuit, def:qc_light_cone_geometry,
        lem:qc_light_cone_geometries}
    \lean{QuantumCircuit.channelCircuitDual_mem_supportedOperators,
        QuantumCircuit.channelCircuitDual_mem_supportedOperators_lightCone,
        QuantumCircuit.BondChannelLayer.dual_mem_supportedOperators,
        QuantumCircuit.ChannelLayer.dual_mem_supportedOperators}
    \leanok
    Let $\Phi$ be a local channel circuit of depth $T$ and let $Z$ act on a
    set of sites $X$. Then $\Phi^\dagger(Z)$ acts on the light cone of radius
    $T$ of $X$; on the ring, on the sites within ring distance $T$ of $X$.
\end{lemma}
\begin{proof}
    \leanok
    By induction on the layers it suffices to treat one layer. The dual of a
    channel on a bond disjoint from $X$ fixes $Z$: its operators commute with
    $Z$, so $\sum_jK_j^\dagger ZK_j=\sum_jK_j^\dagger K_jZ=Z$. The dual of a
    channel on a bond meeting $X$ maps operators acting on the one-step
    neighbourhood of $X$ to such operators, since that set contains the bond.
    The channels of a layer commute, so the dual of the layer is the dual of
    the channels meeting $X$ applied after the dual of the others.
\end{proof}

\begin{lemma}[No connected correlations beyond the light cone of channels]
    \label{lem:qc_channel_light_cone_correlations}
    \uses{lem:qc_channel_light_cone, lem:qc_local_channel_circuit_properties}
    \lean{QuantumCircuit.trace_mul_mul_eq_of_isChannelPreparedInDepth,
        QuantumCircuit.IsBondChannelPreparedInDepth.trace_mul_mul_eq,
        QuantumCircuit.channelCircuitDual_mul,
        QuantumCircuit.channelCircuitDual_mul_of_disjoint_lightCone,
        QuantumCircuit.BondChannelLayer.dual_mul,
        QuantumCircuit.ChannelLayer.dual_mul,
        QuantumCircuit.trace_rectKronecker_mul_mul}
    \leanok
    Let $\rho$ be prepared in depth $T$ by local channels, and let $Z$, $Z'$
    act on sets of sites whose light cones of radius $T$ are disjoint, for
    instance on the ring sets at distance larger than $2T$. Then
    $\Tr(\rho ZZ')
        =\Tr(\rho Z)\Tr(\rho Z')$.
\end{lemma}
\begin{proof}
    \leanok
    The dual of a channel is not multiplicative, but the dual of one layer is
    multiplicative on operators $Z$, $Z'$ acting on sets $X$, $Y$ whose
    one-step neighbourhoods are disjoint. Indeed, no bond of the layer
    meets both $X$ and $Y$. If a channel acts off a set on which an operator
    $B$ acts, then its dual $\Psi$ satisfies $\Psi(AB)=\Psi(A)B$ and
    $\Psi(BA)=B\Psi(A)$. Writing the dual of the layer as the dual of the
    channels meeting $X$ after the dual of the others, and using
    Lemma~\ref{lem:qc_channel_light_cone} for the latter applied to $Z'$,
    one obtains the dual of $ZZ'$ as the product of the duals. By induction
    on the layers, $\Phi^\dagger(ZZ')=\Phi^\dagger(Z)\Phi^\dagger(Z')$, and
    the two factors act on disjoint sets. For a product density
    $\sigma=\bigotimes_i\sigma_i$ of trace one and products
    $\bigotimes_im_i$, $\bigotimes_im'_i$ with disjoint supports,
    $\Tr(\sigma\bigotimes_im_im'_i)
        =\prod_i\Tr(\sigma_im_im'_i)$, and at each site one of
    $m_i$, $m'_i$ is the identity; bilinearity and the duality
    $\Tr(\Phi(\sigma)A)=\Tr(\sigma\Phi^\dagger(A))$
    conclude.
\end{proof}

Ancillas are attached and discarded site by site, so the operations that
change the local dimension act on single sites.

\begin{definition}[Onsite channel]
    \label{def:qc_onsite_channel}
    \lean{Matrix.rectKronecker, Matrix.rectKronecker_one,
        QuantumCircuit.OnsiteChannel,
        QuantumCircuit.OnsiteChannel.map, QuantumCircuit.OnsiteChannel.dual,
        QuantumCircuit.OnsiteChannel.id, QuantumCircuit.OnsiteChannel.attach,
        QuantumCircuit.OnsiteChannel.attachZero,
        QuantumCircuit.OnsiteChannel.attach_siteMap_apply,
        QuantumCircuit.OnsiteChannel.discard,
        QuantumCircuit.OnsiteChannel.discard_siteMap_apply}
    \leanok
    An \emph{onsite channel} from sites of dimension $d$ to sites of dimension
    $e$ is given, at each site $i$, by operators $K_{ij}\colon\C^d\to\C^e$ with
    $\sum_jK_{ij}^\dagger K_{ij}=\mathbb 1$. It is the channel
    $\rho\mapsto\sum_JK_J\rho K_J^\dagger$ on any finite set of sites, such as
    the ring, with
    $K_J=\bigotimes_iK_{iJ(i)}$ for every choice $J$ of one index per site;
    its dual is $A\mapsto\sum_JK_J^\dagger AK_J$. Write
    $\Phi_i(\sigma)=\sum_jK_{ij}\sigma K_{ij}^\dagger$ for its action at site
    $i$. Attaching an ancilla of dimension $a$ in the state
    $\tau=\sum_kw_kw_k^\dagger$ with $\sum_k\|w_k\|^2=1$ is the onsite channel
    $X\mapsto X\otimes\tau$ from dimension $d$ to dimension $da$, and
    discarding it is the onsite channel given by the partial trace over the
    second factor of $\C^d\otimes\C^a$. Ancillas attached to each site and
    local operations on a site and its ancillas are the free operations of
    \cite{Piroli2021LOCC}.
\end{definition}

\begin{lemma}[Onsite channels are product channels]
    \label{lem:qc_onsite_channel_properties}
    \uses{def:qc_onsite_channel}
    \lean{QuantumCircuit.OnsiteChannel.map_isKrausCPTP,
        QuantumCircuit.OnsiteChannel.map_rectKronecker,
        QuantumCircuit.OnsiteChannel.map_finKronecker,
        QuantumCircuit.OnsiteChannel.trace_map_mul,
        QuantumCircuit.OnsiteChannel.discard_siteMap_attach_siteMap}
    \leanok
    An onsite channel $\Phi$ is trace preserving and completely positive, it
    maps $\bigotimes_i\sigma_i$ to $\bigotimes_i\Phi_i(\sigma_i)$, it satisfies
    $\Tr(\Phi(\rho)A)=\Tr(\rho\,\Phi^\dagger(A))$, and discarding an attached
    ancilla gives back the original operator.
\end{lemma}
\begin{proof}
    \leanok
    The operators $K_J$ satisfy
    $\sum_JK_J^\dagger K_J=\bigotimes_i\sum_jK_{ij}^\dagger K_{ij}=\mathbb 1$,
    so the Kraus form gives complete positivity and trace preservation. Since
    $K_J$ is a tensor product, $K_J(\bigotimes_i\sigma_i)K_J^\dagger
        =\bigotimes_iK_{iJ(i)}\sigma_iK_{iJ(i)}^\dagger$, and summing over $J$
    factorizes site by site. The duality is cyclicity of the trace. Finally,
    the partial trace of $X\otimes\tau$ is $\Tr(\tau)X=X$.
\end{proof}

\begin{lemma}[Onsite channels do not spread supports]
    \label{lem:qc_onsite_channel_locality}
    \uses{def:qc_onsite_channel}
    \lean{QuantumCircuit.OnsiteChannel.dual_mem_supportedOperators,
        QuantumCircuit.OnsiteChannel.dual_mul,
        QuantumCircuit.OnsiteChannel.dual_rectKronecker,
        QuantumCircuit.OnsiteChannel.dual_finKronecker,
        QuantumCircuit.rectangularKrausMap_rectKronecker_rectKronecker}
    \leanok
    Let $\Phi$ be an onsite channel. If $Z$ acts on a set of sites $X$, then
    $\Phi^\dagger(Z)$ acts on $X$. If moreover $Z'$ acts on a set of sites
    disjoint from $X$, then $\Phi^\dagger(ZZ')=\Phi^\dagger(Z)\Phi^\dagger(Z')$.
    The dual maps $\bigotimes_im_i$ to $\bigotimes_i\Phi_i^\dagger(m_i)$: for
    any operators $K_{ij}\colon\C^d\to\C^e$, the map
    $\rho\mapsto\sum_JK_J\rho K_J^\dagger$ with $K_J=\bigotimes_iK_{iJ(i)}$
    sends $\bigotimes_im_i$ to $\bigotimes_i\sum_jK_{ij}m_iK_{ij}^\dagger$.
\end{lemma}
\begin{proof}
    \leanok
    The dual maps $\bigotimes_im_i$ to $\bigotimes_i\Phi_i^\dagger(m_i)$, and
    $\Phi_i^\dagger(\mathbb 1)=\sum_jK_{ij}^\dagger K_{ij}=\mathbb 1$. Hence
    the identity factors off $X$ are kept. For products with disjoint
    supports, one of $m_i$, $m'_i$ is the identity at each site, so
    $\Phi_i^\dagger(m_im'_i)=\Phi_i^\dagger(m_i)\Phi_i^\dagger(m'_i)$.
    Linearity concludes.
\end{proof}

\begin{definition}[Local channel conversion]
    \label{def:qc_local_channel_conversion}
    \uses{def:qc_local_channel_circuit, def:qc_onsite_channel}
    \lean{QuantumCircuit.IsLocalChannelProtocol,
        QuantumCircuit.IsLocalChannelConversion}
    \leanok
    A \emph{local channel protocol of depth $T$} is a composition
    $\Psi=\Phi_T\circ L_T\circ\Phi_{T-1}\circ\dots\circ L_1\circ\Phi_0$ of
    channel layers $L_t$ and maps $\Phi_t$, each a composition of finitely
    many onsite channels, which may change the local dimension; $\Phi_0$
    contains at least one onsite channel, and $\Phi_t$ for $t\ge1$ may be
    the identity. A matrix $\rho$ on the chain is
    \emph{converted in depth $T$} into a matrix $\sigma$, possibly of
    another local dimension, if $\sigma=\Psi(\rho)$ for a local channel
    protocol $\Psi$ of depth at most $T$. This is the channel form of the circuits
    $V'=U_\ell V_\ell\cdots U_1V_1U_0$ of \cite{Piroli2021LOCC}, with
    ancillas attached to each site, local operations between the layers and
    the ancillas traced out at the end, without measurements.
\end{definition}

\begin{lemma}[Basic properties of local channel conversions]
    \label{lem:qc_local_channel_conversion_properties}
    \uses{def:qc_local_channel_conversion, def:qc_local_channel_circuit,
        lem:qc_onsite_channel_properties}
    \lean{QuantumCircuit.IsLocalChannelProtocol.isKrausCPTP,
        QuantumCircuit.IsLocalChannelProtocol.channelCircuitMap_comp,
        QuantumCircuit.isLocalChannelConversion_circuit,
        QuantumCircuit.IsChannelPreparedInDepth.exists_isLocalChannelConversion,
        QuantumCircuit.IsLocalChannelConversion.refl,
        QuantumCircuit.IsLocalChannelConversion.mono,
        QuantumCircuit.IsLocalChannelConversion.exists_isKrausCPTP,
        QuantumCircuit.IsLocalChannelConversion.density}
    \leanok
    A local channel protocol is trace preserving and completely positive.
    Attaching ancillas, running a local channel circuit of depth $T$ and
    discarding is a local channel protocol of depth $T$. Every density matrix
    is converted into itself in depth $0$, a conversion in depth $T$ is a
    conversion in every larger depth, a conversion of a density matrix is a
    density matrix, and a density matrix prepared in depth $T$ by local
    channels is converted in depth $T$ from a product density.
\end{lemma}
\begin{proof}
    \leanok
    By induction on the protocol: onsite channels are trace preserving and
    completely positive by Lemma~\ref{lem:qc_onsite_channel_properties},
    channel layers are by construction, and both properties are stable under
    composition. Appending the layers of a circuit of depth $T$ one at a time,
    each onsite channel keeps the depth and each layer raises it by one; the
    attaching and discarding channels are onsite. The identity onsite channel
    gives depth $0$, depth bounds are monotone, and a conversion is realized by
    a channel, which maps a positive semidefinite matrix of trace one to a
    positive semidefinite matrix of trace one. A preparation in depth $T$ by
    local channels starts from a product density and is itself such a
    protocol.
\end{proof}

\begin{lemma}[Composition of conversions]
    \label{lem:qc_local_channel_conversion_trans}
    \uses{def:qc_local_channel_conversion}
    \lean{QuantumCircuit.IsLocalChannelConversion.trans,
        QuantumCircuit.IsLocalChannelProtocol.comp}
    \leanok
    If $\rho$ is converted in depth $T_1$ into $\sigma$ and $\sigma$ is
    converted in depth $T_2$ into $\tau$, then $\rho$ is converted in depth
    $T_1+T_2$ into $\tau$.
\end{lemma}
\begin{proof}
    \leanok
    The composition $\Psi_2\circ\Psi_1$ of local channel protocols of depths
    $T_1$ and $T_2$ is a local channel protocol of depth $T_1+T_2$. Write
    $\Psi_2$ as the sequence of its onsite channels and channel layers, and
    apply them after $\Psi_1$ one at a time: the first onsite channel of
    $\Psi_2$ becomes a further onsite channel in the last group $\Phi_{T_1}$
    of $\Psi_1$, each onsite channel keeps the depth, and each of the $T_2$
    channel layers raises it by one.
\end{proof}

\begin{lemma}[Light cone of a local channel protocol]
    \label{lem:qc_protocol_light_cone}
    \uses{def:qc_local_channel_conversion, lem:qc_onsite_channel_locality,
        lem:qc_channel_light_cone, lem:qc_channel_light_cone_correlations}
    \lean{QuantumCircuit.IsLocalChannelProtocol.exists_dual}
    \leanok
    A local channel protocol $\Psi$ of depth $T$ has a dual $\Psi^\dagger$ with
    $\Tr(\Psi(\rho)A)=\Tr(\rho\,\Psi^\dagger(A))$ that maps operators acting
    on a set of sites $X$ to operators acting on the sites within distance
    $T$ of $X$, and that satisfies
    \begin{align}
        \Psi^\dagger(ZZ') & =\Psi^\dagger(Z)\,\Psi^\dagger(Z')
        \label{eq:qc_conversion_dual_mul}
    \end{align}
    whenever $Z$ and $Z'$ act on sets whose neighbourhoods of radius $T$ are
    disjoint. This is the light cone of a depth-$T$ circuit in the proof of
    \cite[Theorem~1]{Malz2024Preparation}, here with channels and onsite
    operations.
\end{lemma}
\begin{proof}
    \leanok
    By induction on the protocol. For an onsite channel this is
    Lemma~\ref{lem:qc_onsite_channel_locality} with radius $0$; a channel
    layer adds one to the radius, by Lemma~\ref{lem:qc_channel_light_cone}
    for the support and as in the proof of
    Lemma~\ref{lem:qc_channel_light_cone_correlations} for
    (\ref{eq:qc_conversion_dual_mul}), and an onsite channel keeps it.
\end{proof}

\begin{lemma}[No connected correlations after a conversion]
    \label{lem:qc_local_channel_conversion_correlations}
    \uses{def:qc_local_channel_conversion, lem:qc_protocol_light_cone}
    \lean{QuantumCircuit.trace_mul_mul_eq_of_isLocalChannelConversion}
    \leanok
    Let $\sigma$ be converted in depth $T$ from a product
    $\bigotimes_i\sigma_i$ of one-site matrices of trace one, and let $Z$,
    $Z'$ act on sets of sites at distance larger than $2T$. Then
    $\Tr(\sigma ZZ')=\Tr(\sigma Z)\Tr(\sigma Z')$. For pure states prepared
    from a product state by local unitary circuits this is the necessary
    condition for preparation by quantum circuits of \cite{Piroli2021LOCC}.
\end{lemma}
\begin{proof}
    \leanok
    Write $\rho=\bigotimes_i\sigma_i$ and $\sigma=\Psi(\rho)$ for a local
    channel protocol $\Psi$ of depth $T'\le T$, and let $\Psi^\dagger$ be its
    dual from Lemma~\ref{lem:qc_protocol_light_cone}, of radius $T'$.
    If $Z$ and $Z'$ act on sets at distance larger than $2T\ge2T'$, then
    $\Psi^\dagger(Z)$ and $\Psi^\dagger(Z')$ act on disjoint sets, and
    (\ref{eq:qc_conversion_dual_mul}) together with the factorization of the
    product density $\rho$ gives
    \begin{align}
        \Tr(\sigma ZZ') & =\Tr\bigl(\rho\,\Psi^\dagger(ZZ')\bigr)
        =\Tr\bigl(\rho\,\Psi^\dagger(Z)\Psi^\dagger(Z')\bigr) \notag \\
                        & =\Tr\bigl(\rho\,\Psi^\dagger(Z)\bigr)
        \Tr\bigl(\rho\,\Psi^\dagger(Z')\bigr)
        =\Tr(\sigma Z)\Tr(\sigma Z'). \notag
    \end{align}
\end{proof}

\section{Approximate local channel conversions}

The trace-norm approximation of~\cite[paragraph ``Phases of matter'']{Piroli2021LOCC}
can be applied to the local channels of
Definition~\ref{def:qc_local_channel_conversion}. This restricts the allowed
operations: the source additionally permits measurements and classical
feedforward, whereas the channels in this section do not; see
\cite{gap:psc21_local_channel_phase_scope}.

\begin{definition}[Approximate local channel conversion]
    \label{def:qc_approx_local_channel_conversion}
    \uses{def:qc_local_channel_conversion}
    \lean{QuantumCircuit.chainTraceNorm,
        QuantumCircuit.IsApproxLocalChannelConversion}
    \leanok
    A chain operator $\rho$ is approximately converted into $\sigma$ in
    depth $T$ with error $\varepsilon$ if there is a local channel protocol
    $\Psi$ of depth at most $T$ with
    $\lVert\Psi(\rho)-\sigma\rVert_1\le\varepsilon$.
    Here $\lVert A\rVert_1$ is the sum of the singular values of $A$.
\end{definition}

\begin{theorem}[Zero error and monotonicity]
    \label{thm:qc_approx_conversion_properties}
    \uses{def:qc_approx_local_channel_conversion}
    \lean{QuantumCircuit.IsLocalChannelConversion.approx,
        QuantumCircuit.IsApproxLocalChannelConversion.refl,
        QuantumCircuit.IsApproxLocalChannelConversion.mono,
        QuantumCircuit.isApproxLocalChannelConversion_zero_iff}
    \leanok
    Every operator is converted into itself with depth and error zero.
    Increasing the depth or error bound preserves approximate conversion.
    Approximate conversion with error zero is equivalent to exact conversion
    with the same depth bound.
\end{theorem}
\begin{proof}
    \leanok
    Use the identity channel, enlarge the bounds, and use
    $\lVert A\rVert_1=0$ if and only if $A=0$, respectively.
\end{proof}

\begin{theorem}[Composition of approximate conversions]
    \label{thm:qc_approx_conversion_trans}
    \uses{def:qc_approx_local_channel_conversion,
        lem:qc_local_channel_conversion_trans}
    \lean{QuantumCircuit.IsLocalChannelProtocol.chainTraceNorm_map_le,
        QuantumCircuit.IsApproxLocalChannelConversion.trans}
    \leanok
    Let $\rho$ and $\sigma$ be Hermitian. If $\rho$ is approximately converted
    into $\sigma$ with depth $T_1$ and error $\varepsilon_1$, and $\sigma$ into
    $\tau$ with depth $T_2$ and error $\varepsilon_2$, then $\rho$ is
    approximately converted into $\tau$ with depth $T_1+T_2$ and error
    $\varepsilon_1+\varepsilon_2$.
\end{theorem}
\begin{proof}
    \leanok
    The composition $\Psi_2\circ\Psi_1$ has the sum of the two depths.
    The error $\Psi_1(\rho)-\sigma$ is Hermitian. Trace-norm contractivity of
    trace-preserving positive maps and the triangle inequality give
    \begin{align}
      \lVert\Psi_2(\Psi_1(\rho))-\tau\rVert_1
      &\le \lVert\Psi_2(\Psi_1(\rho)-\sigma)\rVert_1
          +\lVert\Psi_2(\sigma)-\tau\rVert_1 \notag\\
      &\le \lVert\Psi_1(\rho)-\sigma\rVert_1
          +\lVert\Psi_2(\sigma)-\tau\rVert_1
      \le\varepsilon_1+\varepsilon_2. \notag
    \end{align}
\end{proof}

\begin{definition}[Asymptotic local channel conversion]
    \label{def:qc_asymptotic_local_channel_conversion}
    \uses{def:qc_approx_local_channel_conversion}
    \lean{QuantumCircuit.IsAsymptoticLocalChannelConversion,
        QuantumCircuit.IsLocalChannelPhaseEquivalent}
    \leanok
    Let $\rho_N$ and $\sigma_N$ be families of operators on chains of $N$
    sites, with fixed local dimensions that may differ between the families.
    The family $\rho_N$ is asymptotically locally convertible into $\sigma_N$
    if there are an integer $p\ge0$, a constant $C\ge0$, and errors
    $\varepsilon_N\to0$ such that, for all sufficiently large $N$, there
    is an approximate conversion with error $\varepsilon_N$ and depth
    $T_N\le C(1+\log(N+1))^p$. Two families are locally phase equivalent
    if they are asymptotically locally convertible in both directions.
\end{definition}

\begin{theorem}[Local channel phase equivalence]
    \label{thm:qc_local_channel_phase_equivalence}
    \uses{def:qc_asymptotic_local_channel_conversion,
        thm:qc_approx_conversion_trans}
    \lean{QuantumCircuit.IsAsymptoticLocalChannelConversion.refl,
        QuantumCircuit.IsAsymptoticLocalChannelConversion.trans,
        QuantumCircuit.IsLocalChannelPhaseEquivalent.refl,
        QuantumCircuit.IsLocalChannelPhaseEquivalent.symm,
        QuantumCircuit.IsLocalChannelPhaseEquivalent.trans,
        QuantumCircuit.IsLocalChannelPhaseEquivalent.equivalence}
    \leanok
    Local channel phase equivalence is an equivalence relation on eventually
    Hermitian families, and in particular on density-matrix families.
\end{theorem}
\begin{proof}
    \leanok
    Identity channels give reflexivity; the definition gives symmetry.
    For transitivity, compose the approximate conversions in each direction.
    The sum of the errors tends to zero. For $b_N=1+\log(N+1)\ge1$, the
    depth satisfies
    \begin{align}
      C_1b_N^{p_1}+C_2b_N^{p_2}
      &\le(C_1+C_2)b_N^{\max(p_1,p_2)}. \notag
    \end{align}
    Thus the depth remains polylogarithmic.
\end{proof}


\section{Fixed physical ports and bounded local registers}

The intersite gates of~\cite[Supplement, pp.~7--8]{Piroli2021LOCC}
act on fixed-dimensional physical qudits. Local ancillas are separate:
free onsite operations may mix the physical qudit with its local ancillas,
but a counted intersite gate does not act on an arbitrarily enlarged pair
of sites. This section derives a resource bound for encoded matching
channels in that convention. It does not identify the unrestricted
local-channel model with the source's phase relation; see
\cite{gap:psc21_physical_port_simulation_scope}.

\begin{definition}[Physical-port channel protocols]
    \label{def:qc_physical_port_protocol}
    \lean{QuantumCircuit.PhysicalPortLayout,
        QuantumCircuit.IsPhysicalPortUnitary,
        QuantumCircuit.IsPhysicalPortUnitary.mem_unitary,
        QuantumCircuit.IsPhysicalPortUnitary.conjTranspose,
        QuantumCircuit.IsPhysicalPortUnitary.one,
        QuantumCircuit.IsPhysicalPortUnitary.onsite_swap,
        QuantumCircuit.IsPhysicalPortProtocol,
        QuantumCircuit.IsPhysicalPortProtocol.isKrausCPTP,
        QuantumCircuit.IsPhysicalPortProtocol.one}
    \leanok
    Every wire carries a $d$-level qudit. A layout assigns its wires to spatial
    sites and selects one physical port at each site. Free operations are
    unitaries supported inside one spatial site, wire permutations preserving
    each spatial site, and normalized onsite Kraus channels. A counted layer
    is an existing nearest-neighbor unitary layer placed on the physical ports,
    with identity on every memory wire. Products add physical depth bounds;
    a larger upper bound remains valid. Reversing a unitary circuit preserves
    its depth. Every resulting channel protocol is CPTP on all input operators.
\end{definition}

\begin{lemma}[Register codes and completed decoding]
    \label{lem:qc_register_code_completion}
    \lean{QuantumCircuit.registerEncoding,
        QuantumCircuit.registerEncoding_isIsometry,
        QuantumCircuit.registerEncoding_isKrausCPTP,
        QuantumCircuit.registerDecoding,
        QuantumCircuit.registerDecoding_isKrausCPTP,
        QuantumCircuit.registerDecoding_encoding,
        QuantumCircuit.registerDecoding_comp_encoding,
        QuantumCircuit.registerDecoding_encoding_reference,
        QuantumCircuit.registerEncodedMap,
        QuantumCircuit.registerEncodedMap_isKrausCPTP,
        QuantumCircuit.registerEncodedMap_encoding,
        QuantumCircuit.le_registerDimension,
        QuantumCircuit.registerBasisEmbedding,
        QuantumCircuit.pairRegisterBasisEmbedding,
        QuantumCircuit.pairRegisterBasisEmbedding_apply,
        QuantumCircuit.registerEncoding_apply,
        QuantumCircuit.registerEncoding_pair}
    \leanok
    An injection $e$ of local basis labels into register words defines an
    isometry $E$. The encoding is $X\mapsto EXE^\dagger$. For a fixed density
    matrix $\rho$, the completed decoder is
    \begin{align}
        D(Y) & =E^\dagger YE+
        \operatorname{tr}[(I-EE^\dagger)Y(I-EE^\dagger)]\rho.
    \end{align}
    It is CPTP and satisfies $D(EXE^\dagger)=X$ for every operator $X$.
    The same identity holds after tensoring with the identity on any finite
    reference. Encoding the two endpoints separately gives the tensor product
    of their local encoding matrices. If $m\le B$ and $d\ge2$, then
    $k=\lceil\log_d B\rceil$ qudits suffice. A channel $\Phi$ extends to the
    full input register by $\widehat\Phi=\mathcal E_{\rm out}\Phi D_{\rm in}$;
    this extension is CPTP and agrees with $\Phi$ on encoded inputs.
\end{lemma}
\begin{proof}
    \leanok
    The inclusion has orthonormal columns. Complete the compression outside its
    range using the existing support-completion channel; the complementary
    term vanishes on $EXE^\dagger$. Matrix multiplication gives the exact left
    inverse, and tensoring preserves composition. The paired basis inclusion
    factors entrywise into two local inclusions. The ceiling-logarithm bound
    is $B\le d^k$.
\end{proof}

\begin{lemma}[Parallel physical SWAP layers]
    \label{lem:qc_port_matching_swaps}
    \uses{def:qc_physical_port_protocol}
    \lean{QuantumCircuit.permOpHom,
        QuantumCircuit.permOp_one,
        QuantumCircuit.permOp_mul,
        QuantumCircuit.permOp_conjTranspose,
        QuantumCircuit.PortMatching.perm,
        QuantumCircuit.PortMatching.perm_apply_of_notMem,
        QuantumCircuit.PortMatching.perm_apply_of_mem,
        QuantumCircuit.PortMatching.perm_apply_left,
        QuantumCircuit.PortMatching.perm_apply_right,
        QuantumCircuit.PortMatching.perm_involutive,
        QuantumCircuit.PortMatching.perm_symm,
        QuantumCircuit.PortMatching.layer,
        QuantumCircuit.PortMatching.layer_op,
        QuantumCircuit.PortMatching.layer_op_one_site}
    \leanok
    A matching of ring bonds defines an involutive site permutation: each
    selected bond exchanges its endpoints and every other site is fixed.
    The actual layer containing these SWAP gates has exactly the associated
    permutation matrix. Permutation matrices preserve products and inverses.
    The self-loop of a one-site ring acts as the identity.
\end{lemma}
\begin{proof}
    \leanok
    Swaps on disjoint bonds commute and fix each other's endpoints. Their
    finite product therefore has the stated action. Applying the permutation
    matrix homomorphism identifies that product with the layer operator.
\end{proof}

\begin{theorem}[Selected-memory routing at unchanged physical depth]
    \label{thm:qc_port_register_routing}
    \uses{def:qc_physical_port_protocol, lem:qc_port_matching_swaps}
    \lean{QuantumCircuit.PhysicalPortLayout.selectionPermutation,
        QuantumCircuit.PhysicalPortLayout.selectionPermutation_site,
        QuantumCircuit.PhysicalPortLayout.selectionPermutation_port,
        QuantumCircuit.PhysicalPortLayout.selectionPermutation_symm,
        QuantumCircuit.IsPhysicalPortUnitary.selected_layer,
        QuantumCircuit.embedOp_permOp_viaEmbedding,
        QuantumCircuit.IsPhysicalPortUnitary.selected_permutation_layer,
        QuantumCircuit.registerPermutation,
        QuantumCircuit.registerPermutation_apply,
        QuantumCircuit.registerPermutation_apply_of_notMem,
        QuantumCircuit.registerPermutation_involutive,
        QuantumCircuit.IsPhysicalPortUnitary.register_matching}
    \leanok
    Select one memory wire at each site. Swap it with the site's physical port,
    apply a physical layer, and undo the onsite swaps. This acts exactly as the
    layer on the selected wires and as identity on all other wires, at physical
    depth one. Repeating this for $k$ disjoint register digits implements the
    same matching permutation on each digit at depth at most $k$, independently
    of the number of selected bonds. All unselected wires are restored on the
    whole operator space, without assuming they were initialized or uncorrelated.
\end{theorem}
\begin{proof}
    \leanok
    The selection permutation preserves each spatial site, hence is free.
    Conjugating the port layer by it changes only the wire placement. Extend
    each digit's site permutation by identity outside that digit's selected
    wires. Disjoint ranges ensure the finite product has the desired action
    on each digit; every factor has a one-layer implementation.
\end{proof}

\begin{lemma}[Data and scratch placement]
    \label{lem:qc_port_data_scratch}
    \uses{thm:qc_port_register_routing}
    \lean{QuantumCircuit.PortRegisters.layout,
        QuantumCircuit.PortRegisters.wire,
        QuantumCircuit.PortRegisters.site_wire,
        QuantumCircuit.PortRegisters.data,
        QuantumCircuit.PortRegisters.scratch,
        QuantumCircuit.PortRegisters.site_data,
        QuantumCircuit.PortRegisters.site_scratch,
        QuantumCircuit.PortRegisters.data_ne_scratch,
        QuantumCircuit.PortRegisters.selection,
        QuantumCircuit.PortRegisters.selection_apply,
        QuantumCircuit.PortRegisters.site_selection,
        QuantumCircuit.PortRegisters.selection_ranges_disjoint,
        QuantumCircuit.PortRegisters.pairData,
        QuantumCircuit.PortRegisters.pairData_left,
        QuantumCircuit.PortRegisters.pairData_right,
        QuantumCircuit.PortMatching.right_notMem,
        QuantumCircuit.PortRegisters.routingPermutation,
        QuantumCircuit.PortRegisters.data_notMem_selection,
        QuantumCircuit.PortRegisters.routingPermutation_data_left,
        QuantumCircuit.PortRegisters.routingPermutation_data_of_notMem,
        QuantumCircuit.PortRegisters.routingPermutation_pairData_site,
        QuantumCircuit.PortRegisters.routingPermutation_isPhysicalPortUnitary,
        QuantumCircuit.PortRegisters.pairData_isPhysicalPortProtocol}
    \leanok
    Use $1+2k$ wires per spatial site: one physical port, $k$ data qudits and
    $k$ scratch qudits. On a matching, select a data digit at each left endpoint
    and a scratch digit at each right endpoint. The common routing moves the
    left data register into the right scratch register and leaves the right
    data register fixed. Thus both data registers of each nontrivial selected
    bond lie at its right spatial site after routing. Reversing the routing
    returns their outputs to their original sites and restores all other wires.
\end{lemma}
\begin{proof}
    \leanok
    A right endpoint of a nontrivial matching bond cannot be another left
    endpoint. Data and scratch slots are disjoint, and different digit selections
    have disjoint ranges. The action formulas of
    Theorem~\ref{thm:qc_port_register_routing} give the claimed placement.
\end{proof}

\begin{theorem}[One shared routing for a channel family]
    \label{thm:qc_port_shared_channel_routing}
    \uses{def:qc_physical_port_protocol, thm:qc_port_register_routing}
    \lean{QuantumCircuit.sum_conjTranspose_embedOp_mul,
        QuantumCircuit.routingMatrixEquiv,
        QuantumCircuit.routingMatrixEquiv_apply,
        QuantumCircuit.routingMatrixEquiv_symm_apply,
        QuantumCircuit.routingChannelConj,
        QuantumCircuit.routingChannelConj_apply,
        QuantumCircuit.routingChannelConj_embeddedKraus,
        QuantumCircuit.IsPhysicalPortProtocol.of_routed_map,
        QuantumCircuit.IsPhysicalPortProtocol.of_register_routing,
        QuantumCircuit.IsPhysicalPortProtocol.of_routed_family}
    \leanok
    Suppose an actual wire routing of physical depth $T$ moves every channel
    in a commuting family into a single spatial site, possibly a different
    site for each member. Execute all those conjugated channels locally,
    then reverse the routing. The result is exactly the original family
    product, with physical depth at most $2T$. The number of family members
    does not multiply this depth bound. The identities hold for every input
    operator, including arbitrary correlations with unused registers.
\end{theorem}
\begin{proof}
    \leanok
    Wire reindexing is a linear equivalence on matrices. Its conjugation on
    linear maps is multiplicative, so it preserves the commuting family
    product. Placing Kraus operators preserves their completeness relation.
    Each routed map is a zero-depth onsite channel; only the one forward and
    one backward routing contribute to the depth.
\end{proof}

\begin{theorem}[Uniform bounded-register matching simulation]
    \label{thm:qc_bounded_physical_port_simulation}
    \uses{lem:qc_register_code_completion, lem:qc_port_data_scratch,
        thm:qc_port_shared_channel_routing}
    \lean{QuantumCircuit.PortRegisters.pairChannel,
        QuantumCircuit.PortRegisters.pairChannel_commute,
        QuantumCircuit.PortRegisters.matchingChannel,
        QuantumCircuit.PortRegisters.matchingChannel_isPhysicalPortProtocol,
        QuantumCircuit.boundedPairRegisterCode,
        QuantumCircuit.exists_bounded_matching_simulation,
        QuantumCircuit.IsPhysicalPortProtocol.list_prod_uniform}
    \leanok
    Fix $d\ge2$, $B$, and a local input dimension $m\le B$, and choose a
    fallback density matrix on a pair of local inputs. Set
    $k=\lceil\log_d B\rceil$, independently of the chain length.
    For every ring of $N\ge2$ sites and every matching of CPTP pair channels,
    there are normalized Kraus maps on the two $k$-qudit data registers that
    agree with the prescribed code extensions on the entire register space,
    and intertwine exactly with the original pair channels on encoded inputs.
    Their actual placed matching channel has a physical-port protocol of depth
    at most $2k$. A product of $L$ such certified layers has depth at most $2kL$.
    For $B=1$ the register width and this bound are zero.
\end{theorem}
\begin{proof}
    \leanok
    Encode the two endpoints independently, then complete each pair channel
    outside its joint code. The resulting map is CPTP, so choose its normalized
    finite Kraus representation. The common data-to-scratch routing has depth
    $k$ and puts every pair channel onsite. Apply
    Theorem~\ref{thm:qc_port_shared_channel_routing} once to the complete
    matching. Protocol composition gives the final linear-in-$L$ bound.
    The off-code joint decoder need not factor onsite; it is only executed
    after colocation. This theorem is not an assertion about protocols with
    an unrestricted intermediate dimension depending on $N$.
\end{proof}

\begin{lemma}[Product onsite coding on arbitrary chain inputs]
    \label{lem:qc_product_register_encoding}
    \uses{lem:qc_register_code_completion}
    \lean{QuantumCircuit.OnsiteChannel.ofSiteMaps,
        QuantumCircuit.OnsiteChannel.ofSiteMaps_siteMap,
        QuantumCircuit.OnsiteChannel.comp,
        QuantumCircuit.OnsiteChannel.comp_map,
        QuantumCircuit.OnsiteChannel.map_eq_id_of_siteMap_eq_id,
        QuantumCircuit.OnsiteChannel.encodeRegisters,
        QuantumCircuit.OnsiteChannel.encodeRegisters_siteMap,
        QuantumCircuit.OnsiteChannel.encodeRegisters_map,
        QuantumCircuit.OnsiteChannel.decodeRegisters,
        QuantumCircuit.OnsiteChannel.decodeRegisters_encodeRegisters,
        QuantumCircuit.OnsiteChannel.decodeRegisters_encodeRegisters_reference,
        QuantumCircuit.OnsiteChannel.encoded,
        QuantumCircuit.OnsiteChannel.encoded_map_encodeRegisters}
    \leanok
    Finite CPTP maps at individual sites assemble into an onsite channel,
    and assembly preserves composition. Product isometric encodings have
    product completed local decoders whose composition is the identity on
    all chain operators, also with any finite reference. Extending an existing
    onsite channel between site-consistent input and output codes intertwines
    with its original chain map, including entangled chain inputs.
\end{lemma}
\begin{proof}
    \leanok
    Choose normalized local Kraus families. On product operators, all formulas
    follow site by site; these products span the entire chain matrix algebra.
    Use the exact local decoding identity before tensoring with a reference.
\end{proof}

\begin{lemma}[Representation-independent placement of register channels]
    \label{lem:qc_register_channel_lift}
    \uses{lem:qc_placed_operators}
    \lean{QuantumCircuit.registerConfigurationSplit,
        QuantumCircuit.registerConfigurationSplit_apply,
        QuantumCircuit.registerConfigurationSplit_symm_comp,
        QuantumCircuit.registerMatrixSplit,
        QuantumCircuit.registerMatrixSplit_conjTranspose,
        QuantumCircuit.registerMatrixSplit_embedOp,
        QuantumCircuit.registerChannelLift,
        QuantumCircuit.registerMatrixSplit_channelLift,
        QuantumCircuit.registerChannelLift_id,
        QuantumCircuit.registerChannelLift_comp,
        QuantumCircuit.registerChannelLift_kraus,
        QuantumCircuit.rectangularKrausMap_embedOp_congr}
    \leanok
    Splitting a configuration into a selected register and its complementary
    wires identifies a placed operator with $A\otimes I$. Under the same
    reindexing, placing a Kraus family gives precisely its local channel
    tensored with the identity on the complement. Consequently the placed
    channel is independent of its chosen Kraus representation, and the
    placement identity holds on arbitrary correlated inputs.
\end{lemma}
\begin{proof}
    \leanok
    Use the explicit configuration equivalence and the matrix reindexing
    algebra equivalence. Check tensoring on matrix units, then use linear
    extensionality. Applying the split to each placed Kraus operator proves
    the full channel identity.
\end{proof}

\begin{theorem}[Intertwining of placed encoded channels]
    \label{thm:qc_encoded_channel_placement}
    \uses{lem:qc_product_register_encoding, lem:qc_register_channel_lift}
    \lean{QuantumCircuit.reindex_rectKronecker_registerSplit,
        QuantumCircuit.registerMatrixSplit_encodeRegisters,
        QuantumCircuit.registerChannelLift_encodeRegisters}
    \leanok
    Let $E_S$ be the product code of a selected register and $E$ the product
    code of the complete chain, possibly with different local input and output
    dimensions. If $\widehat\Phi E_S=E_S\Phi$ as linear maps on operators,
    then the placed channels satisfy
    $\operatorname{Lift}(\widehat\Phi)E=E\operatorname{Lift}(\Phi)$
    on every chain operator. The remaining sites may carry arbitrary
    correlations with the selected register.
\end{theorem}
\begin{proof}
    \leanok
    The rectangular global encoding matrix splits as the tensor product of
    its selected and complementary encoding matrices. Check the resulting
    tensor intertwining on matrix units, then use linear extensionality and
    the representation-independent placement identity.
\end{proof}

%% ==================================================
%% Section: Lieb--Robinson propagation for finite-range Hamiltonians on graphs
%% Sources: Hastings--Koma, arXiv:math-ph/0507008, Appendix A, (A.12)--(A.16);
%% Nachtergaele--Ogata--Sims, arXiv:math-ph/0603064, Section 2; OpenAI, "A
%% two-dimensional area law from a global spectral gap" (September 24, 2026,
%% openai/math@adc7f1241b42e322a6451854ab7e4b4c146bf78a),
%% build/sections/03-quasilocal.tex, lines 14--41 and 52--116 (Lemma 4.1,
%% label lem:quasilocal-lr). The results below are stated for a general
%% finite graph and a general interaction family; the induced square-lattice
%% statement of Lemma 4.1 is a separate entry.
%% ==================================================

\section{Lieb--Robinson propagation on finite graphs}
\label{sec:qc_lieb_robinson}

Let $I$ be a finite set of sites of local dimension $q$, and write
$\mathcal A_X$ for the operators on $(\C^q)^{\otimes I}$ acting on
$X\subseteq I$. All norms are Hilbert-space operator norms. For an operator
$H$ put $\tau_t(A)=e^{itH}Ae^{-itH}$. The dynamics generated by a sum of
local Hermitian terms moves support at a bounded speed, up to exponentially
small tails \cite{Hastings2006Spectral,Nachtergaele2006LiebRobinson}. The
estimate below measures distance by the path metric of an arbitrary simple
graph on $I$, so that holes and separate connected components cost nothing:
sites in different components are at infinite distance, and there the
commutator vanishes exactly.

\begin{definition}[Restricted commutator norm]
    \label{aux:restricted-commutator-norm}
    \lean{QuantumCircuit.heisenbergCommutatorNorm,
        QuantumCircuit.heisenbergEvolution}\leanok
    For operators $H,B$ and a set of sites $X$, define
    \[
        C_B(X,t)=\sup\{\|[\tau_t(A),B]\|: A\in\mathcal A_X,\ \|A\|\leq1\}.
    \]
\end{definition}

\begin{lemma}[Evaluation of the restricted commutator norm]
    \label{aux:restricted-commutator-evaluation}
    \lean{QuantumCircuit.norm_heisenberg_commutator_le_heisenbergCommutatorNorm}\leanok
    \uses{aux:restricted-commutator-norm}
    For $A\in\mathcal A_X$, one has
    $\|[\tau_t(A),B]\|\leq C_B(X,t)\|A\|$.
\end{lemma}
\begin{proof}\leanok
    Apply the defining operator-norm bound to $A$.
\end{proof}

\begin{lemma}[Reversal of time]
    \label{aux:restricted-commutator-negative-time}
    \lean{QuantumCircuit.heisenbergCommutatorNorm_neg_time}\leanok
    \uses{aux:restricted-commutator-norm}
    The restricted commutator norm for $H$ at time $-t$ equals that for
    $-H$ at time $t$.
\end{lemma}
\begin{proof}\leanok
    Both Heisenberg evolutions are $A\mapsto e^{-itH}Ae^{itH}$.
\end{proof}

\begin{lemma}[Continuity in time]
    \label{aux:restricted-commutator-continuity}
    \lean{QuantumCircuit.continuous_heisenbergCommutatorNorm}\leanok
    \uses{aux:restricted-commutator-norm}
    For fixed $H,B,X$, the function $t\mapsto C_B(X,t)$ is continuous.
\end{lemma}
\begin{proof}\leanok
    Matrix exponentiation and multiplication are continuous; the norm of the
    resulting linear map on $\mathcal A_X$ is continuous.
\end{proof}

\begin{lemma}[Initial commutator bound]
    \label{aux:restricted-initial-commutator-bound}
    \lean{QuantumCircuit.heisenbergCommutatorNorm_zero_le}\leanok
    \uses{aux:restricted-commutator-norm}
    For arbitrary $H,B,X$, one has $C_B(X,0)\leq2\|B\|$.
\end{lemma}
\begin{proof}\leanok
    Use $\|AB-BA\|\leq2\|A\|\|B\|$.
\end{proof}

\begin{lemma}[Vanishing initial commutator]
    \label{aux:restricted-initial-commutator-vanishing}
    \lean{QuantumCircuit.heisenbergCommutatorNorm_zero_of_disjoint}\leanok
    \uses{aux:restricted-commutator-norm}
    If $B\in\mathcal A_Y$ and $X\cap Y=\varnothing$, then $C_B(X,0)=0$.
\end{lemma}
\begin{proof}\leanok
    Operators on disjoint tensor factors commute.
\end{proof}

\begin{lemma}[Local commutator recursion]
    \label{aux:restricted-commutator-recursion}
    \lean{QuantumCircuit.heisenbergCommutatorNorm_le_integral}\leanok
    \uses{aux:restricted-commutator-norm}
    Let $H=\sum_j h_j$ be a finite sum, where each $h_j$ is Hermitian and
    belongs to $\mathcal A_{T_j}$. For arbitrary $B,X$ and $t\geq0$,
    \[
        C_B(X,t)\leq C_B(X,0)
        +2\sum_{j:X\cap T_j\neq\varnothing}\|h_j\|
        \int_0^t C_B(T_j,s)\,ds.
    \]
    No bound on the ground energy or positivity assumption on $h_j$ is
    required, and several terms may share a support. This is the recursion of
    Hastings--Koma, Appendix A, (A.12)--(A.14).
\end{lemma}
\begin{proof}\leanok
    Separate the terms disjoint from $X$. Their sum commutes with every operator
    in $\mathcal A_X$. Differentiate in the corresponding interaction picture,
    apply the fundamental theorem of calculus and invariance of operator norm
    under unitary conjugation, and then take the restricted operator norm.
\end{proof}

\begin{lemma}[Weighted integral comparison]
    \label{aux:weighted-volterra-comparison}
    \lean{finite_volterra_le_weight_mul_exp,
        finite_volterra_zero_le_weight_mul_exp_sub_one,
        finite_rectangular_volterra_zero_le_weight_mul_exp_sub_one,
        le_mul_exp_of_le_const_add_integral}\leanok
    Let $f_i\geq0$, for $i$ in a finite set, be continuous on $[0,T]$ and satisfy
    $f_i(t)\leq b_i+\sum_j a_{ij}\int_0^tf_j(s)\,ds$ with $a_{ij}\geq0$.
    If positive weights $w_i$ satisfy $\sum_ja_{ij}w_j\leq vw_i$ and
    $b_i\leq cw_i$, then $f_i(t)\leq cw_ie^{vt}$ for $0\leq t\leq T$. If moreover
    $b_i=0$ and $v>0$, then $f_i(t)\leq cw_i(e^{vt}-1)$.
\end{lemma}
\begin{proof}\leanok
    The quantity $\max_if_i(t)/w_i$ satisfies the scalar inequality
    $g(t)\leq c+v\int_0^tg$, and Gr\"onwall's inequality gives $g(t)\leq ce^{vt}$.
    For the refinement, insert this bound into the integral once more.
\end{proof}

\begin{lemma}[Exact locality on interaction-closed regions]
    \label{lem:qc_lr_closed_region}
    \lean{QuantumCircuit.commute_heisenbergEvolution_of_closed}\leanok
    Let $H=\sum_jh_j$ with $h_j\in\mathcal A_{T_j}$, and let $C$ be a set of sites
    such that each $T_j$ is contained in $C$ or disjoint from $C$. If
    $A\in\mathcal A_C$ and $B\in\mathcal A_D$ with $C\cap D=\varnothing$, then
    $[\tau_t(A),B]=0$ for every $t$.
\end{lemma}
\begin{proof}\leanok
    Write $H=H_C+H'$, with $H_C$ the sum of the terms inside $C$. The two parts
    commute, and $H'$ commutes with $\mathcal A_C$, so
    $\tau_t(A)=e^{itH_C}Ae^{-itH_C}$. Every factor commutes with $B$.
\end{proof}

\paragraph{Finite-range interactions on a graph.}
Let $G$ be a simple graph on $I$ with path metric $d_G$, taking the value
$\infty$ between different connected components. Let $H=\sum_kh_k$ be a
finite sum of Hermitian terms with $h_k\in\mathcal A_{X_k}$, and choose an
anchor $a_k\in X_k$. Assume that every support has $d_G$-diameter at most
$R$ and at most $v_R$ sites, and that the norms of the terms whose support
contains any given site sum to at most $b_0$.

\begin{theorem}[Graph-distance propagation]
    \label{thm:qc_graph_lieb_robinson}
    \lean{QuantumCircuit.norm_heisenberg_commutator_le_graphDistance,
        QuantumCircuit.exists_graph_lieb_robinson}\leanok
    \uses{aux:restricted-commutator-recursion,aux:weighted-volterra-comparison,
        aux:restricted-commutator-evaluation,aux:restricted-commutator-negative-time,
        aux:restricted-initial-commutator-bound,
        aux:restricted-initial-commutator-vanishing,
        aux:restricted-commutator-continuity,thm:qc_graph_lieb_robinson_zero}
    In this setting let $A\in\mathcal A_{X_i}$, let $y$ be a site, and let
    $B\in\mathcal A_{\{y\}}$. For every $\mu\geq0$, every integer
    $0\leq n\leq d_G(a_i,y)$ and every real $t$,
    \[
        \|[\tau_t(A),B]\|\leq 2e^{\mu R}\|A\|\|B\|
        \exp\bigl(2v_Rb_0e^{2\mu R}|t|-\mu n\bigr).
    \]
    Consequently there are positive constants $C,v,c$, depending only on
    $R,v_R,b_0$, such that
    $\|[\tau_t(A),B]\|\leq C\|A\|\|B\|e^{v|t|-cd_G(a_i,y)}$ whenever
    $d_G(a_i,y)<\infty$, and $[\tau_t(A),B]=0$ otherwise. The constants do not
    depend on the number of sites, on $q$, or on the graph.
\end{theorem}
\begin{proof}\leanok
    Let $F_k(t)=C_B(X_k,t)$ for $t\geq0$. The local commutator recursion gives
    $F_k(t)\leq F_k(0)+\sum_ja_{kj}\int_0^tF_j$, with $a_{kj}=2\|h_j\|$ when
    $X_j$ meets $X_k$ and $a_{kj}=0$ otherwise. Put
    $D(x)=\min\{n,d_G(x,y)\}$ and $w_k=e^{-\mu D(a_k)}$. If $X_j$ meets
    $X_k$, then $d_G(a_k,a_j)\leq2R$, so $w_j\leq e^{2\mu R}w_k$; the norms of the
    terms meeting $X_k$ sum to at most $|X_k|b_0\leq v_Rb_0$. Hence
    $\sum_ja_{kj}w_j\leq2v_Rb_0e^{2\mu R}w_k$. The initial value $F_k(0)$
    vanishes unless $y\in X_k$, and then $D(a_k)\leq R$ and
    $F_k(0)\leq2\|B\|\leq2\|B\|e^{\mu R}w_k$. The weighted integral comparison
    gives $F_i(t)\leq2\|B\|e^{\mu R}w_ie^{2v_Rb_0e^{2\mu R}t}$ with
    $w_i=e^{-\mu n}$. Replace $H$ by $-H$ for negative times and multiply by
    $\|A\|$. For the constants take $\mu=1$; the infinite-distance case is
    Lemma~\ref{thm:qc_graph_lieb_robinson_zero}.
\end{proof}

\begin{lemma}[Vanishing between connected components]
    \label{thm:qc_graph_lieb_robinson_zero}
    \lean{QuantumCircuit.heisenberg_commutator_eq_zero_of_edist_eq_top}\leanok
    \uses{lem:qc_lr_closed_region}
    In the setting of Theorem~\ref{thm:qc_graph_lieb_robinson}, if
    $d_G(a_i,y)=\infty$ then $[\tau_t(A),B]=0$ for every $t$.
\end{lemma}
\begin{proof}\leanok
    Let $C$ be the connected component of $a_i$. A support of finite diameter
    lies in one component, so it is contained in $C$ or disjoint from it; also
    $X_i\subseteq C$ and $y\notin C$. Apply
    Lemma~\ref{lem:qc_lr_closed_region}.
\end{proof}

\paragraph{Normalized partial traces.}
For $K\subseteq I$ let
$E_K(B)=q^{-|I\setminus K|}\operatorname{Tr}_{I\setminus K}(B)\otimes\mathbb 1_{I\setminus K}$,
with every tensor factor at its original site.

\begin{definition}[Normalized partial-trace expectation]
    \label{def:qc_site_expectation}
    \lean{QuantumCircuit.siteExpectation, QuantumCircuit.traceOutside,
        QuantumCircuit.siteExpectationLM}\leanok
    The map $E_K$ on the operators of $(\C^q)^{\otimes I}$ defined above.
\end{definition}

\begin{lemma}[Weyl averaging formula]
    \label{lem:qc_site_expectation_average}
    \lean{QuantumCircuit.siteExpectation_eq_average,
        QuantumCircuit.sum_localWeyl_mul_star, QuantumCircuit.outsideWeyl}\leanok
    \uses{def:qc_site_expectation}
    Let $q\geq1$, and for $p=(a,b)$ let $W(p)=X^aZ^b$ be the Weyl operators on
    $\C^q$. Then $E_K(B)$ is the uniform average of $W_pBW_p^\dagger$ over all
    products $W_p=\bigotimes_{x\notin K}W(p_x)$ of Weyl operators outside $K$.
\end{lemma}
\begin{proof}\leanok
    The Weyl operators satisfy
    $\sum_pW(p)_{ii'}\overline{W(p)_{jj'}}=q\,\delta_{ij}\delta_{i'j'}$. The
    corresponding sum over the products factorizes over the sites, and the
    resulting Kronecker deltas reproduce the partial-trace formula.
\end{proof}

\begin{lemma}[Properties of the expectation]
    \label{lem:qc_site_expectation_properties}
    \lean{QuantumCircuit.siteExpectation_mem_supportedOperators,
        QuantumCircuit.siteExpectation_of_mem_supportedOperators,
        QuantumCircuit.siteExpectation_one, QuantumCircuit.norm_siteExpectation_le,
        QuantumCircuit.siteExpectationLM_isKrausCPTP}\leanok
    \uses{lem:qc_site_expectation_average}
    For $q\geq1$, the map $E_K$ takes values in $\mathcal A_K$, fixes every
    operator in $\mathcal A_K$, is unital, completely positive and trace
    preserving, and is a contraction in operator norm.
\end{lemma}
\begin{proof}\leanok
    The values act on $K$ by the partial-trace formula. The averaging formula
    is a Kraus decomposition by rescaled unitaries. Operators acting on $K$
    commute with every $W_p$. A uniform average of unitary conjugations is a
    contraction.
\end{proof}

\begin{lemma}[Localization error through on-site commutators]
    \label{lem:qc_site_expectation_commutator}
    \lean{QuantumCircuit.norm_sub_siteExpectation_le}\leanok
    \uses{lem:qc_site_expectation_average}
    If $\|[B,U]\|\leq\varepsilon_y$ for every site $y\notin K$ and every unitary
    $U\in\mathcal A_{\{y\}}$, then $\|B-E_K(B)\|\leq\sum_{y\notin K}\varepsilon_y$.
\end{lemma}
\begin{proof}\leanok
    By the averaging formula, $B-E_K(B)$ is the average of
    $(BW_p-W_pB)W_p^\dagger$. Write $W_p$ as a product of on-site unitaries and
    use $\|[B,UV]\|\leq\|[B,U]\|+\|[B,V]\|$ for unitaries $U,V$.
\end{proof}

\begin{theorem}[Localization onto graph balls]
    \label{thm:qc_graph_ball_localization}
    \lean{QuantumCircuit.norm_heisenberg_sub_siteExpectation_graphBall_le,
        QuantumCircuit.graphBall, QuantumCircuit.sum_pow_mul_exp_neg_le}\leanok
    \uses{thm:qc_graph_lieb_robinson,thm:qc_graph_lieb_robinson_zero,
        lem:qc_site_expectation_commutator,lem:qc_site_expectation_properties}
    In the setting of Theorem~\ref{thm:qc_graph_lieb_robinson}, with $q\geq1$,
    assume that at most $K_g(d+1)^k$ sites lie at distance exactly $d$ from any
    site. Let $N_l(a)=\{x:d_G(a,x)\leq l\}$. Then
    \[
        \|\tau_t(A)-E_{N_l(a_i)}(\tau_t(A))\|
        \leq\|A\|\min\bigl\{2,\;Ce^{v|t|-l/2}\bigr\},
    \]
    with $v=2v_Rb_0e^{2R}$ and $C=2e^RK_g\,2^kk!\,e^{1/2}/(1-e^{-1/2})$.
\end{theorem}
\begin{proof}\leanok
    The bound $2\|A\|$ follows from contractivity. For the other bound, apply the
    commutator lemma with $\varepsilon_y=0$ when $d_G(a_i,y)=\infty$ and
    $\varepsilon_y=2e^R\|A\|e^{v|t|-d_G(a_i,y)}$ otherwise, using graph-distance
    propagation with $\mu=1$. Group the sites outside the ball by their distance
    $d>l$, and use $(d+1)^k\leq2^kk!\,e^{(d+1)/2}$ and a geometric series.
\end{proof}

% Source: OpenAI2026AreaLaw, 09-amplification.tex, lines 101--139,
% openai/math@adc7f1241b42e322a6451854ab7e4b4c146bf78a.
% This leaf covers physical-site oscillation and spectator-preserving averaging.
% It does not assert the channel-shell estimate, the one-event increment,
% or propagation by retained Poisson clocks.
% SiteOscillation is checked at 3d99bf2: native build, strict elaboration,
% six examples and fifteen axiom guards pass.
% SpectatorSiteExpectation is checked at 030de2afb512bae217d5f7ee15654ca8b443b58a:
% native build, strict elaboration, seven examples and eight axiom guards pass.

\section{Physical-site oscillation with a spectator}
\label{sec:qc_site_oscillation}

Let $I$ and $J$ be finite sets, let $q\in\N$, and put
$\mathcal H=(\C^q)^{\otimes I}$ and $\mathcal K=\C^J$.
The spectator space $\mathcal K$ has arbitrary finite dimension, including
zero. Write $\mathcal A_X$ for the physical operators acting on
$X\subseteq I$, as in Definition~\ref{def:qc_local_circuit}.
Operators $B,C$ on $\mathcal H\otimes\mathcal K$ are arbitrary complex
matrices. Their norm is the operator norm induced by the Euclidean norm:
$\|B\|=\sup_{\|\psi\|_2\le1}\|B\psi\|_2$.
For physical operators $U$, write $\widehat U=U\otimes\mathbb 1_{\mathcal K}$,
and use $[A,B]=AB-BA$.

\subsection{The on-site unitary supremum}

\begin{definition}[Oscillation at a physical site]
    \label{def:qc_site_oscillation}
    \uses{def:qc_local_circuit}
    \lean{QuantumCircuit.siteOscillation}
    \leanok
    For $y\in I$, define
    \begin{align}
        d_y(B)&=\sup\bigl\{\|[\widehat U,B]\|:
            U\in\mathcal A_{\{y\}},\ U^\dagger U=UU^\dagger=\mathbb 1\bigr\}.
        \label{eq:qc_osc_supremum}
    \end{align}
    This is the on-site unitary supremum of
    \cite{OpenAI2026AreaLaw}, with every test unitary acting as
    the identity on the spectator.
\end{definition}

\begin{theorem}[Bounds for the oscillation supremum]
    \label{thm:qc_site_oscillation_bounds}
    \uses{def:qc_site_oscillation}
    \lean{QuantumCircuit.siteOscillation_bddAbove,
        QuantumCircuit.siteOscillation_nonneg,
        QuantumCircuit.norm_commutator_le_siteOscillation,
        QuantumCircuit.norm_commutator_le_siteOscillation_right,
        QuantumCircuit.siteOscillation_le,
        QuantumCircuit.siteOscillation_le_two_mul_norm}
    \leanok
    The set in~(\ref{eq:qc_osc_supremum}) contains zero and is bounded
    above by $2\|B\|$. In particular, its supremum is finite and
    $0\le d_y(B)\le2\|B\|$. Every unitary $U\in\mathcal A_{\{y\}}$
    satisfies $\|[\widehat U,B]\|=\|[B,\widehat U]\|\le d_y(B)$.
    If these commutator norms are all at most a real number $r$, then
    $d_y(B)\le r$.
\end{theorem}
\begin{proof}
    \leanok
    \uses{def:qc_site_oscillation}
    The identity is an admissible unitary and its commutator is zero.
    For each admissible $U$, the operator $\widehat U$ is unitary, so
    \begin{align}
        \|[\widehat U,B]\|
        &\le\|\widehat UB\|+\|B\widehat U\|=2\|B\|.
        \label{eq:qc_osc_commutator_bound}
    \end{align}
    The assertions follow from the upper- and lower-bound properties
    of the supremum. Reversing a commutator changes only its sign.
\end{proof}

\begin{theorem}[Subadditivity and absolute homogeneity]
    \label{thm:qc_site_oscillation_seminorm}
    \uses{def:qc_site_oscillation}
    \lean{QuantumCircuit.siteOscillation_add_le,
        QuantumCircuit.siteOscillation_smul,
        QuantumCircuit.siteOscillation_neg,
        QuantumCircuit.siteOscillation_sub_le,
        QuantumCircuit.siteOscillation_sum_le}
    \leanok
    For every $c\in\C$,
    \begin{align}
        d_y(B+C)&\le d_y(B)+d_y(C),
        \label{eq:qc_osc_subadditivity}\\
        d_y(cB)&=|c|d_y(B).
        \label{eq:qc_osc_homogeneity}
    \end{align}
    Consequently $d_y(-B)=d_y(B)$ and
    $d_y(B-C)\le d_y(B)+d_y(C)$.
    For any finite set $F$ and any family $(B_k)_{k\in F}$,
    $d_y(\sum_{k\in F}B_k)\le\sum_{k\in F}d_y(B_k)$.
\end{theorem}
\begin{proof}
    \leanok
    \uses{thm:qc_site_oscillation_bounds,thm:qc_site_oscillation_commutant}
    The identities
    $[\widehat U,B+C]=[\widehat U,B]+[\widehat U,C]$ and
    $[\widehat U,cB]=c[\widehat U,B]$, followed by the norm inequalities
    and the supremum bound, give~(\ref{eq:qc_osc_subadditivity}) and
    $d_y(cB)\le|c|d_y(B)$.
    If $c\ne0$, apply the latter inequality to $c^{-1}$ and $cB$ to
    obtain the reverse bound in~(\ref{eq:qc_osc_homogeneity}).
    If $c=0$, use $d_y(0)=0$.
    Take $c=-1$ for negation and write $B-C=B+(-C)$ for subtraction.
    For the finite sum, distribute the commutator and use
    \begin{align}
        \Bigl\|\Bigl[\widehat U,\sum_{k\in F}B_k\Bigr]\Bigr\|
        &\le\sum_{k\in F}\|[\widehat U,B_k]\|
        \le\sum_{k\in F}d_y(B_k).
    \end{align}
    Taking the supremum over $U$ proves the claim.
\end{proof}

\begin{theorem}[Vanishing on the physical commutant]
    \label{thm:qc_site_oscillation_commutant}
    \uses{def:qc_site_oscillation}
    \lean{QuantumCircuit.siteOscillation_eq_zero_of_commute,
        QuantumCircuit.siteOscillation_zero,
        QuantumCircuit.siteOscillation_one_kronecker}
    \leanok
    If $[\widehat U,B]=0$ for every $U\in\mathcal A_{\{y\}}$, then
    $d_y(B)=0$. In particular, $d_y(0)=0$, and every operator $S$ on
    the spectator satisfies $d_y(\mathbb 1_{\mathcal H}\otimes S)=0$.
\end{theorem}
\begin{proof}
    \leanok
    \uses{thm:qc_site_oscillation_bounds}
    The commutation hypothesis bounds every norm in the supremum by
    zero, whereas $d_y(B)\ge0$. For a spectator operator,
    $(U\otimes\mathbb 1)(\mathbb 1\otimes S)=U\otimes S
    =(\mathbb 1\otimes S)(U\otimes\mathbb 1)$.
\end{proof}

\subsection{Averaging physical sites}
\label{sec:qc_spectator_expectation}

For $K\subseteq I$, let $E_K$ be the normalized partial-trace expectation
of Definition~\ref{def:qc_site_expectation}. The set $K$ consists of the
retained physical sites. Thus averaging over a set $Q\subseteq I$
means using $K=I\setminus Q$.

\begin{definition}[Site expectation with a spectator]
    \label{def:qc_spectator_expectation}
    \uses{def:qc_site_expectation}
    \lean{QuantumCircuit.spectatorSiteExpectation}
    \leanok
    For $a,b\in J$, let $B_{ab}$ be the physical matrix whose entries are
    $(B_{ab})_{\sigma,\tau}=B_{(\sigma,a),(\tau,b)}$.
    Define $E_K^{\mathcal K}$ by applying $E_K$ to each such block:
    \begin{align}
        \bigl(E_K^{\mathcal K}(B)\bigr)_{(\sigma,a),(\tau,b)}
        &=\bigl(E_K(B_{ab})\bigr)_{\sigma,\tau}.
        \label{eq:qc_spectator_blocks}
    \end{align}
    In particular, when $q>0$, set $Q=I\setminus K$ and write
    $u\sqcup\rho$ for the configuration equal to $u$ on $K$ and to
    $\rho$ on $Q$. Then
    \begin{align}
        \bigl(E_K^{\mathcal K}(B)\bigr)_{(\sigma,a),(\tau,b)}
        &=q^{-|Q|}\mathbf1_{\{\sigma|_Q=\tau|_Q\}}
          \sum_{\rho\in\{0,\ldots,q-1\}^{Q}}
          B_{(\sigma|_K\sqcup\rho,a),(\tau|_K\sqcup\rho,b)}.
        \label{eq:qc_spectator_trace}
    \end{align}
    Both spectator indices remain unchanged.
\end{definition}

\begin{theorem}[Preservation of spectator factors]
    \label{thm:qc_spectator_factors}
    \uses{def:qc_spectator_expectation}
    \lean{QuantumCircuit.spectatorSiteExpectation_kronecker,
        QuantumCircuit.spectatorSiteExpectation_one_kronecker,
        QuantumCircuit.spectatorSiteExpectation_univ}
    \leanok
    Suppose $q>0$. For physical operators $A$ and spectator operators $S$,
    \begin{align}
        E_K^{\mathcal K}(A\otimes S)&=E_K(A)\otimes S,
        \label{eq:qc_spectator_tensor}\\
        E_K^{\mathcal K}(\mathbb 1\otimes S)&=\mathbb 1\otimes S.
        \label{eq:qc_spectator_fixed}
    \end{align}
    Retaining all physical sites gives $E_I^{\mathcal K}(B)=B$.
\end{theorem}
\begin{proof}
    \leanok
    \uses{def:qc_spectator_expectation,def:qc_site_expectation,
        lem:qc_site_expectation_properties}
    The $(a,b)$ block of $A\otimes S$ is $S_{ab}A$. Linearity gives
    $E_K(S_{ab}A)=S_{ab}E_K(A)$, proving~(\ref{eq:qc_spectator_tensor}).
    The expectation fixes the identity, which yields
    (\ref{eq:qc_spectator_fixed}). Every physical matrix is supported
    on $I$, so $E_I(B_{ab})=B_{ab}$ for every block.
\end{proof}

\begin{theorem}[Commutation on the averaged sites]
    \label{thm:qc_spectator_commutant}
    \uses{def:qc_spectator_expectation}
    \lean{QuantumCircuit.spectatorSiteExpectation_commute}
    \leanok
    For every $U\in\mathcal A_{I\setminus K}$,
    $[E_K^{\mathcal K}(B),\widehat U]=0$.
\end{theorem}
\begin{proof}
    \leanok
    \uses{def:qc_spectator_expectation,def:qc_site_expectation}
    Each block $E_K(B_{ab})$ is supported on $K$ by its partial-trace
    formula. Operators on disjoint sets of physical sites commute, so
    $E_K(B_{ab})U=UE_K(B_{ab})$. Matrix multiplication by
    $\widehat U$ acts separately on each spectator block, giving the
    claimed equality.
\end{proof}

\begin{theorem}[Amplified Weyl averaging formula]
    \label{thm:qc_spectator_average}
    \uses{def:qc_spectator_expectation,lem:qc_site_expectation_average}
    \lean{QuantumCircuit.spectatorSiteExpectation_eq_average}
    \leanok
    Suppose $q>0$. For $p\in(\Z/q\Z\times\Z/q\Z)^I$, let
    $W_K(p)$ be the product of the on-site Weyl operators $W(p_x)$ at
    $x\notin K$, with the identity at $x\in K$.
    Then
    \begin{align}
        E_K^{\mathcal K}(B)
        &=q^{-2|I|}\sum_{p\in(\Z/q\Z\times\Z/q\Z)^I}
          \widehat{W_K(p)}B\widehat{W_K(p)}^{\dagger}.
        \label{eq:qc_spectator_average}
    \end{align}
    Labels at retained sites are included in the sum; each choice of
    Weyl operators outside $K$ consequently occurs $q^{2|K|}$ times.
\end{theorem}
\begin{proof}
    \leanok
    \uses{def:qc_spectator_expectation,lem:qc_site_expectation_average}
    For every physical $U$, the $(a,b)$ block of
    $\widehat UB\widehat U^\dagger$ is $UB_{ab}U^\dagger$.
    Apply the physical Weyl averaging formula to each $B_{ab}$ in
    (\ref{eq:qc_spectator_blocks}).
\end{proof}

\begin{theorem}[Contraction with an arbitrary spectator]
    \label{thm:qc_spectator_contraction}
    \uses{def:qc_spectator_expectation}
    \lean{QuantumCircuit.norm_spectatorSiteExpectation_le}
    \leanok
    If $q>0$, then $\|E_K^{\mathcal K}(B)\|\le\|B\|$.
\end{theorem}
\begin{proof}
    \leanok
    \uses{thm:qc_spectator_average}
    Every $\widehat{W_K(p)}$ is unitary, so each summand in
    (\ref{eq:qc_spectator_average}) has norm $\|B\|$.
    There are $q^{2|I|}$ labels. The triangle inequality and the
    normalization $q^{-2|I|}q^{2|I|}=1$ give the result.
\end{proof}

\begin{theorem}[Averaging error from on-site commutator bounds]
    \label{thm:qc_spectator_error}
    \uses{def:qc_spectator_expectation}
    \lean{QuantumCircuit.norm_sub_spectatorSiteExpectation_le}
    \leanok
    Suppose $q>0$, and let $\varepsilon:I\to\R$ satisfy
    $\|[B,\widehat U]\|\le\varepsilon_y$ for each $y\notin K$ and
    every unitary $U\in\mathcal A_{\{y\}}$. Then
    \begin{align}
        \|B-E_K^{\mathcal K}(B)\|
        &\le\sum_{y\notin K}\varepsilon_y.
        \label{eq:qc_spectator_error}
    \end{align}
\end{theorem}
\begin{proof}
    \leanok
    \uses{thm:qc_spectator_average}
    For unitaries $V,W$ on $\mathcal H\otimes\mathcal K$,
    $[B,VW]=[B,V]W+V[B,W]$ and hence
    $\|[B,VW]\|\le\|[B,V]\|+\|[B,W]\|$.
    Induction over the sites outside $K$ bounds the commutator with
    $\widehat{W_K(p)}$ by $\sum_{y\notin K}\varepsilon_y$.
    With $V=\widehat{W_K(p)}$, unitarity gives
    \begin{align}
        B-VBV^\dagger&=[B,V]V^\dagger,
        \label{eq:qc_spectator_difference}\\
        \|B-VBV^\dagger\|&=\|[B,V]\|
            \le\sum_{y\notin K}\varepsilon_y.
        \label{eq:qc_spectator_difference_bound}
    \end{align}
    Average~(\ref{eq:qc_spectator_difference_bound}) over all Weyl labels
    using~(\ref{eq:qc_spectator_average}) and
    $q^{-2|I|}\sum_p B=B$ to obtain~(\ref{eq:qc_spectator_error}).
\end{proof}

For $q>0$, take $K=I\setminus Q$ and $\varepsilon_y=d_y(B)$ in
Theorem~\ref{thm:qc_spectator_error}, with
Theorem~\ref{thm:qc_site_oscillation_bounds}, to obtain the physical-site
averaging estimate of \cite{OpenAI2026AreaLaw}:
\begin{align}
    \|B-E_{I\setminus Q}^{\mathcal K}(B)\|
    &\le\sum_{y\in Q}d_y(B).
    \label{eq:qc_osc_averaging_error}
\end{align}
By Theorem~\ref{thm:qc_spectator_commutant}, the averaged operator
commutes with $U\otimes\mathbb 1_{\mathcal K}$ for every
$U\in\mathcal A_Q$. These bounds contain no factor depending on
$\dim\mathcal K$.

